\subsection{Weighted Particle--Mesh Boltzmann--Poisson Module}
\label{sec:weighted-vlasov-pic}

\noindent\textbf{Maturity and validation.}
The task-specific weighted particle-in-cell (WPIC) runners are supported
production implementations, but that status is not a blanket validation of a
new physical regime.  Every scientific run still needs marker-count, radial-
mesh, timestep, collision-cadence, collision-locality, and random-realization
convergence.  Constant-cross-section material self-interacting dark matter
(SIDM) runs have the strongest benchmark path.  Velocity-dependent SIDM,
unequal-weight adaptation, and subhalo capture remain research
extensions until the problem-specific long-time and population-control tests
described below have passed.

\begin{figure}[htbp]
\centering
\textbf{Physics flow: weighted particle--mesh module}
\par\smallskip
\begin{tikzpicture}[node distance=0.22cm]
\node[physicsflow] (wp1) {Sample persistent weighted markers from the initial spherical phase space distribution};
\node[physicsflow,below=of wp1] (wp2) {Deposit every active marker's represented mass onto the radial gravity mesh};
\node[physicsflow,below=of wp2] (wp3) {Solve the spherical Poisson equation to obtain the gravitational field};
\node[physicsflow,below=of wp3] (wp4) {Apply a half gravitational kick and move the markers along their collisionless trajectories};
\node[physicsflow,below=of wp4] (wp5) {Deposit the moved markers, recompute the endpoint field, and apply the second half kick};
\node[physicsflow,below=of wp5] (wp6) {When a collision event is due, group nearby active markers and sample the physical encounter rate};
\node[physicsflow,below=of wp6] (wp7) {Scatter two complete objects elastically, or merge a captured pair into one complete remnant};
\node[physicsflow,below=of wp7] (wp8) {Update escaped, delayed, and internal-energy state; reconstruct profiles, write diagnostics or a restart, and continue};
\draw[physicsarrow] (wp1) -- (wp2);
\draw[physicsarrow] (wp2) -- (wp3);
\draw[physicsarrow] (wp3) -- (wp4);
\draw[physicsarrow] (wp4) -- (wp5);
\draw[physicsarrow] (wp5) -- (wp6);
\draw[physicsarrow] (wp6) -- (wp7);
\draw[physicsarrow] (wp7) -- (wp8);
\end{tikzpicture}
\end{figure}

\subsubsection{Purpose and state variables}
\label{sec:wpic-continuum}

The weighted particle--mesh module represents Vlasov--Poisson and
Boltzmann--Poisson evolution with moving phase space markers.  Let
\(f(t,r,w,\eta)\) be the mass distribution function, \(r\) physical radius,
\(w\) radial peculiar velocity, and \(\eta=L^2/2\), where \(L\) is the
magnitude of specific angular momentum.  Integrating over the spatial angles
and the tangential velocity azimuth gives
\[
 \dd M=8\pi^2 f\,\dd r\,\dd w\,\dd\eta,
 \qquad \rho(t,r)=\frac{2\pi}{r^2}\int f\,\dd w\,\dd\eta .
                                                               \tag{PIC0a}
\]
Thus \(f\) has units of mass per spatial volume per velocity volume.  With
Hubble rate \(H(t)\), the Boltzmann--Poisson equations are
\[
 \partial_t f+(Hr+w)\partial_r f+
 \left(\frac{2\eta}{r^3}-\Phi_r-Hw\right)\partial_w f=C[f],
                                                               \tag{PIC0b}
\]
Here \(\rho_{\rm src}\) is the gravitational source density: total density
for an isolated static halo, or the prescribed background-subtracted source
for the expanding module. The enclosed source mass and force follow from
Poisson's equation.

\[
 \frac1{r^2}\partial_r(r^2\Phi_r)=4\pi G\rho_{\rm src},
 \qquad \Phi_r=\frac{\partial\Phi}{\partial r}.          \tag{PIC0c}
\]
Here \(G\) is Newton's constant and \(\Phi\) is the peculiar gravitational
potential.  Static isolated calculations set \(H=0\) and use the full halo
density as \(\rho_{\rm src}\).  In an expanding total-density calculation,
\(\rho_{\rm src}=\rho-\bar\rho\), where \(\bar\rho\) is the homogeneous
background density; an already background-subtracted input enters directly.
An imposed central mass and the unresolved interior mass enter through the
inner enclosed-mass boundary condition.  The expanding background and its
exact orbital update are derived in Section~\ref{sec:wpic-expanding-kdk}.

The Vlasov--Poisson equations follow by setting \(C[f]=0\).  In that limit
\(f\) remains constant along the characteristics of (PIC0b).  The particle
method follows those trajectories and deposits their positive masses to
solve (PIC0c).  Local collision sampling supplies \(C[f]\) when encounters
are enabled.  The subhalo operator also evolves each object's mass
and internal energy, so its state includes those additional properties.

The code samples each occupied part of the distribution with many markers.
In the source, \code{WeightedPICMarker} is the record that stores one marker's
complete persistent physical and bookkeeping state:
\[
 (m_p,r_p,w_p,\eta_p,E_{{\rm int},p},
   \sigma_{{\rm int},p},I_p,s_p),
 \qquad \eta_p=\frac{L_p^2}{2},
\]
where \(m_p\) is represented mass, \(r_p\) is areal radius, \(w_p\) is radial
velocity, \(L_p\) is angular momentum magnitude, and
\(\eta_p=L_p^2/2\) is the angular momentum coordinate used by the solver.
Furthermore, \(E_{{\rm int},p}\) is unresolved internal energy,
\(\sigma_{{\rm int},p}\) is its one-dimensional internal velocity dispersion,
\(I_p\) is a persistent integer identifier, and \(s_p\) records whether the
marker is active, delayed inside the unresolved central region, or escaped.
A collisionless orbit update changes \(r_p\) and \(w_p\) while
keeping \(m_p\), \(\eta_p\), \(E_{{\rm int},p}\), and
\(\sigma_{{\rm int},p}\) fixed.  Elastic collisions change \(w_p\) and
\(\eta_p\), but leave both complete marker masses unchanged.  In the Plummer
subhalo operator, capture replaces two whole objects by one remnant of mass
\(m_i+m_j\) and increases its \(E_{\rm int}\).  Both the material SIDM
operator and the Plummer encounter operator accept unequal positive marker
masses.  The field solve uses those same masses on the radial mesh, so the
calculation alternates between continuous marker orbits and mesh-based
gravity without reverting to an equal-mass approximation.

This chapter is restricted to the Newtonian Vlasov--Poisson and
Boltzmann--Poisson WPIC code.  The separate Einstein--Vlasov WPIC
implementation is documented in
Section~\ref{sec:weighted-einstein-vlasov-pic}; it is collisionless and is not
a substitute for the Newtonian collision paths described here.

This representation is useful when a phase space cell contains a sharp orbit
boundary.  Several markers can occupy different positions inside that cell
and follow their individual characteristics.  Orbit integration and field
deposition are implemented in \path{cpu/include/ev/weighted_vlasov_pic.hpp}.
Local elastic scattering is in
\path{cpu/include/ev/weighted_vlasov_pic_collisions.hpp}.  The conservative bridge
back to cell-average \(f(r,w,\eta)\) is in
\path{cpu/include/ev/weighted_vlasov_pic_phase_remap.hpp}.  The two-clock event
controller is in
\path{cpu/include/ev/weighted_vlasov_pic_event_scheduler.hpp}.  The complete
transport--collision--reconstruction event is assembled in
\path{cpu/include/ev/weighted_pic_boltzmann_event.hpp}.
Unequal-mass elastic scattering and inelastic capture for resolved Plummer
subhalo encounters are in
\path{cpu/include/ev/weighted_vlasov_pic_subhalo_collisions.hpp}.

\subsubsection{Step 1: turn cell mass into marker mass}

The spherical phase space measure is
\[
 \dd\Gamma=8\pi^2\,\dd r\,\dd w\,\dd\eta.
\]
Therefore a cell average \(f_{ijk}\) contains the physical mass
\[
 M_{ijk}=8\pi^2f_{ijk}\Delta r_i\Delta w_j\Delta\eta_k. \tag{PIC1}
\]
VENUS first constructs every \(M_{ijk}\) and its cumulative sum.  A set of
evenly spaced mass strata selects source cells from that cumulative sum.  This
gives approximately equal marker masses while retaining the exact total:
\[
 \sum_p m_p=\sum_{ijk}M_{ijk}.                            \tag{PIC2}
\]
The phase-grid initializer chooses a common marker mass from the nearest floating-point values
around the exact quotient.  The same selected mass is assigned to every
marker.  Equation (PIC2) is then satisfied to floating-point accuracy while
an initially quiet equal-mass sample is obtained.  Equality is an initializer
choice, not a collision restriction: restart or adaptive marker populations
may contain unequal positive masses.

Let \(N_{\rm occupied}\) be the number of source phase space cells with
positive represented mass, \(N_{\rm per\ cell}\) the requested nominal number
of markers for each such cell, and \(N_{\rm budget}\) the global upper limit on
the total marker count.  The nominal marker count is
\[
 N_p=\min(N_{\rm occupied}N_{\rm per\ cell},N_{\rm budget}). \tag{PIC3}
\]
The defaults are 32 markers per occupied cell and a global limit of one
million.  Selecting cells through mass strata lets a very small tail remain in
the represented total without assigning one compulsory marker to every tiny
tail cell.

The central source code is:
\codesource{WPIC marker initializer}
{\path{cpu/include/ev/weighted_vlasov_pic.hpp}}
{Converts positive phase space cell masses into a deterministic weighted-marker
 realization with the same represented total mass}
\begin{lstlisting}[language=C++]
cell_mass = 8.0 * PI * PI * f[q]
          * dr[i] * dw[j] * deta[k];

selected = systematic_sources(cells, marker_count,
                              total_mass, offset);
marker.mass = total_mass / marker_count;
\end{lstlisting}


\paragraph{Analytic radial halo initialization.}
The power-law halo task constructs its particles directly in
\code{initialize_markers} in \path{cpu/src/drivers/expanding_wpic.cpp}.
It assigns a specified number of quiet velocity pairs to each occupied
radial source cell. If cell $j$ contains mass $\Delta M_j$ and is sampled
by $N_j$ markers, each marker in that cell has mass
\begin{equation}
 m_p=\frac{\Delta M_j}{N_j}.
 \label{eq:analytic-radial-marker-mass}
\end{equation}
The uniform central source is sampled by the same rule. A common final
normalization factor makes the represented resolved mass equal to its
specified total. The radial cell masses differ, so a fixed marker count
per radial cell produces unequal masses across the halo. Optional angular importance strata divide
each source mass among the selected angular intervals before assigning
its quiet pairs. The literal assignment is
\codesource{Masses in the analytic radial halo sampler}
{\path{cpu/src/drivers/expanding_wpic.cpp}, \code{initialize_markers}, lines 1027--1032}
{Divides the source-shell mass in an angular stratum among its two markers per quiet pair}
\begin{lstlisting}[language=C++]
                const double marker_mass =
                    shell_mass * represented_mass_fraction
                    / static_cast<double>(2 * stratum_pairs);
\end{lstlisting}
Here \code{represented_mass_fraction} is the fraction of the source-shell
mass assigned to the stratum, and \code{stratum_pairs} is its number of
quiet pairs. The ordinary sampler has one angular stratum with fraction
one. In collisionless relaxation, these masses are quadrature weights
for the gravitational source. In the subhalo encounter model,
each is also the physical mass of one subhalo and enters its pair cross
section. Increasing the number of such markers at fixed source mass
therefore changes the physical subhalo population. Its mass distribution
must be specified along with the halo density and velocity distribution.

\subsubsection{Step 2: choose positions inside each source cell}

For a piecewise-constant cell, VENUS uses shifted low-discrepancy coordinates
in \(r,w,\eta\).  They spread the markers across the cell more evenly than
independent random draws.  The shift is generated from the marker identifier,
source-cell identifier, seed, and coordinate number.  Repeating a run with the
same seed therefore produces the same markers.

An analytic equilibrium supplies more information than a cell average.  In
that case the density-aware initializer proposes continuous points inside the
cell and accepts them with probability
\[
 P_{\rm accept}=\frac{f(r,w,\eta)}{f_{\rm bound}},        \tag{PIC4}
\]
where \(f_{\rm bound}\) is a verified upper bound for that cell.  A uniform
number \(z\) in \((0,1)\) is drawn from the counter-based sequence, and the
point is accepted when
\[
 z f_{\rm bound}\leq f(r,w,\eta).                        \tag{PIC5}
\]
Equations (PIC4)--(PIC5) sample the continuous distribution inside a cell.
They are especially important when the energy support cuts through an
\(\eta\) cell.

For a distribution symmetric in radial velocity, markers are emitted in
adjacent pairs with
\[
 (m,r,w,\eta),\qquad(m,r,-w,\eta).                       \tag{PIC6}
\]
Their initial radial momentum cancels exactly.  The code checks the cell
masses and velocity faces before enabling this quiet pairing.

\subsubsection{Step 3: deposit marker mass on the spherical mesh}

Gravity needs a radial density.  Define the volume coordinate
\[
 q=\frac{r^3}{3},\qquad
 \Delta V_i=4\pi(q_{i+1/2}-q_{i-1/2}).                  \tag{PIC7}
\]
Between neighboring shell centers \(q_i\leq q_p\leq q_{i+1}\), cloud-in-cell
deposition uses
\[
 a_{pi}=\frac{q_{i+1}-q_p}{q_{i+1}-q_i},\qquad
 a_{p,i+1}=1-a_{pi}.                                    \tag{PIC8}
\]
The two weights are nonnegative and add to one.  Near a physical boundary,
the in-domain weights are normalized again so that their sum remains one.

The expanding origin-parity path needs one additional construction.  Its
markers remain active while crossing the unresolved physical region
\(0\leq r_p<r_{\min}\), rather than reflecting at the innermost field face
\(r_{\min}\).  Define
\[
 q_f=\frac{r_{\min}^3}{3},\qquad
 \Delta q_0=2(q_0-q_f),\qquad
 W_q=\min(\Delta q_0,q_f),                               \tag{PIC8a}
\]
where \(q_f\) is the volume coordinate of the inner face, \(q_0\) is the
center of the first resolved shell, \(\Delta q_0\) is that shell's width in
\(q\), and \(W_q\) is the inner handoff width.  The minimum with \(q_f\)
keeps the handoff inside the physical range \(q\geq0\).  For a marker with
\(q_p=r_p^3/3<q_f\), VENUS forms
\[
 s_p=\operatorname{clip}\!\left[
       \frac{q_p-(q_f-W_q)}{W_q},0,1\right],
 \qquad f_p=s_p^2(3-2s_p),                               \tag{PIC8b}
\]
where \(\operatorname{clip}(x,0,1)\) limits \(x\) to the closed interval
\([0,1]\), \(s_p\) is the dimensionless position through the handoff, and
\(f_p\) is the fraction of marker \(p\)'s cloud assigned to the resolved
mesh.  The ordinary normalized CIC or TSC face stencil
\(a^{\rm face}_{pi}\), evaluated at \(r_{\min}\), is frozen below the face
and partitioned as
\[
 \widetilde a_{pi}=f_p a^{\rm face}_{pi},\qquad
 h_p=1-f_p,\qquad
 \sum_i\widetilde a_{pi}+h_p=1.                          \tag{PIC8c}
\]
Here \(\widetilde a_{pi}\) is the weight delivered to resolved shell \(i\)
and \(h_p\) is the complementary weight delivered to the unresolved inner
gravity carrier.  For \(q_p\geq q_f\), the ordinary position-dependent
stencil is used and \(h_p=0\); for \(q_p\leq q_f-W_q\), \(f_p=0\) and the
whole marker is carried inside.  Because \(f_p\) and its first derivative
match continuously at both ends, a marker crossing \(r_{\min}\) changes the
resolved and central source smoothly.

VENUS deposits extensive quantities first:
\[
 \begin{aligned}
 M_i&=\sum_p \widetilde a_{pi}m_p,&
 P_{r,i}&=\sum_p \widetilde a_{pi}m_pw_p,\\
 R_{2,i}&=\sum_p \widetilde a_{pi}m_pw_p^2,&
 T_{2,i}&=\sum_p \widetilde a_{pi}m_p\frac{2\eta_p}{r_p^2}.
 \end{aligned}                                           \tag{PIC9}
\]
Here \(i\) labels a resolved radial shell and the sum over \(p\) runs over
active markers.  Thus \(M_i\) is deposited represented mass,
\(P_{r,i}\) is deposited radial momentum, \(R_{2,i}\) is the mass-weighted
radial second moment (the deposited sum of \(w_p^2\)), and \(T_{2,i}\) is the
mass-weighted tangential second moment (the deposited sum of the squared
tangential speed \(2\eta_p/r_p^2\)).
In an ordinary reflecting calculation,
\(\widetilde a_{pi}=a_{pi}\).  In the origin-parity handoff it has the
meaning in equation (PIC8c).
It then divides by \(\Delta V_i\) to form physical shell averages.  Since
\(\sum_i\widetilde a_{pi}+h_p=1\), summing the deposited shell masses and
the unresolved carrier mass gives
\[
 \sum_i M_i+M_{\rm hid}
 =\sum_p m_p,\qquad M_{\rm hid}=\sum_p h_pm_p.           \tag{PIC10}
\]
Thus each marker contributes exactly \(f_pm_p\) to the boundary cloud and
\((1-f_p)m_p\) to the inner carrier: no represented mass is created or
lost.  The carrier is a mass-only unresolved source for gravity; for markers
that are wholly resolved, the usual partition-of-unity proof also applies to
every raw moment in equation (PIC9).  The Poisson solve includes
\(M_{\rm hid}\) as central enclosed mass. Equations (PIC8a)--(PIC8c)
transfer the represented mass continuously between the boundary cloud and
the unresolved central source as the marker crosses \(r_{\min}\).
The marker retains its physical mass and phase space coordinates.

The current Poisson routine represents shell mass as
\(4\pi r_i^2\rho_i\Delta r_i\).  The adapter therefore passes
\[
 \rho_i^{\rm field}=\frac{M_i}{4\pi r_i^2\Delta r_i}.   \tag{PIC11}
\]
Equation (PIC11) makes the mass seen by the Poisson solver exactly equal to
the shell mass \(M_i\); the same field solve adds \(M_{\rm hid}\) separately
as the unresolved central source, so its total represented source remains the
right-hand side of equation (PIC10).

Below the first physical mesh face, \(r_f=r_{\min}\), the Poisson solver
continues the unresolved source as a uniform sphere. Let \(M_f\) be the
enclosed source mass at that face, \(M_c\) a prescribed central point mass,
and \(\Phi_f=\Phi(r_f)\). For expanding \code{total_density} input,
\(M_f=M_c+M_{\rm hid}-4\pi\bar\rho r_f^3/3\), where
\(\bar\rho\) is the homogeneous background density. With
\code{density_contrast} input, the represented mass is already the
source contrast and \(M_f=M_c+M_{\rm hid}\). The continuation is
\[
 \Phi_r(r)=\frac{GM_c}{r^2}
          +\frac{G(M_f-M_c)}{r_f^3}r,\qquad 0<r<r_f,       \tag{PIC11a}
\]
\[
 \Phi(r)=\Phi_f+GM_c\left(\frac{1}{r_f}-\frac{1}{r}\right)
       +\frac{G(M_f-M_c)}{2r_f^3}(r^2-r_f^2).             \tag{PIC11b}
\]
Integrating the uniform density gives an enclosed mass proportional to
\(r^3\), which produces the linear force in (PIC11a). Integrating that
force from \(r_f\) gives (PIC11b). Both quantities join the resolved field
continuously at the face. An empty central cavity with \(M_f=M_c=0\)
therefore has zero force even when the first resolved shell contains matter.
The code samples the face mass and potential before applying this continuation:
\codesource{Gravitational field inside the first mesh face}
{\path{cpu/include/ev/vlasov_poisson.hpp}; \path{gpu/src/wpic_vlasov_poisson.cu}}
{Continue the enclosed source as a uniform central sphere, separating a prescribed point mass}
\begin{lstlisting}[language=C++]
const double r0 = face_resolved ? field.radial_edge.front() : r.front();
const double phi0 = face_resolved ? field.phi_edge.front() : field.phi.front();
const double face_force = face_resolved
    ? p.vp_gravity * field.source_mass_edge.front() / (r0 * r0)
    : field.phi_r.front();
\end{lstlisting}
The unresolved sphere supplies a force approximation inside \(r_f\).
Scientific profiles at nearby radii require comparison with a smaller
\(r_f\) and sufficient radial and particle resolution.

\subsubsection{Step 4: solve gravity and move the markers}

The deposited mass is sent through the ordinary spherical Poisson solver.
Write
\[
 g(t,r)=-\partial_r\Phi(t,r),
\]
where \(\Phi\) is the Newtonian gravitational potential and \(g\) is the
signed radial gravitational acceleration.  For positive enclosed mass,
\(g<0\), so gravity points inward.  Each active marker then follows
\[
 \dot r_p=w_p,\qquad
 \dot w_p=\frac{2\eta_p}{r_p^3}+g(t,r_p)
          =\frac{2\eta_p}{r_p^3}-\Phi_r(t,r_p).          \tag{PIC12}
\]
The library exposes both a generic frozen-field orbit adapter and a
self-consistent endpoint-field KDK update.  The generic adapter calls the same
VENUS characteristic machinery used by the grid solver, including ordinary
integration and selective Levi--Civita retry.  It is used by the frozen-field
event module.

The current static/expanding production runner uses a self-consistent KDK
step:
\begin{enumerate}
\item deposit markers at \(t^n\) and solve \(\Phi^n\);
\item apply one half gravitational kick and an exact central-force drift;
\item deposit the drifted markers at \(t^{n+1}\) and solve \(\Phi^{n+1}\);
\item apply the endpoint half kick.
\end{enumerate}
The drift formula and expanding variables are derived in
Section~\ref{sec:wpic-expanding-kdk}.

Marker states distinguish active, delayed-inner, and escaped mass.  A restart
stores every marker, its state, weight, \(\eta\), persistent identifier,
unresolved internal energy, delayed-boundary data, both random samplers, both
elapsed event-clock ages, time, and step. The self-consistent KDK path
reconstructs its radial field from the stored markers. The static Boltzmann
event path also stores the field held between potential updates.

\subsubsection{Step 5: form local collision shells}

Elastic SIDM scattering is local in position.  At the start of a collision
interval, VENUS sorts the active markers by radius and divides the sorted list
into shells.  The nominal shell contains 4096 markers, following the SphDSMC
construction used for the Frank benchmark.  The two shell faces are placed
halfway in the volume coordinate \(q=r^3/3\) between the last marker on one
side and the first marker on the other side.  This gives
\[
 V_s=\frac{4\pi}{3}\left(r_{s,+}^3-r_{s,-}^3\right),
 \qquad
 n_s=\frac{N_s}{V_s},\qquad
 n_{{\rm partner},s}=\frac{N_s-1}{V_s},\qquad
 \rho_s=\frac{\sum_{i\in s}m_i}{V_s}.                   \tag{PIC13}
\]
Here \(N_s\) is the number of complete markers in shell \(s\), \(m_i\) is the
represented mass of marker \(i\), \(n_s\) is the marker number density, and
\(\rho_s\) is mass density.  A chosen marker cannot collide with itself, so
its finite-shell density of available partners is
\(n_{{\rm partner},s}=(N_s-1)/V_s\).  The familiar continuum replacement
\(n_{{\rm partner},s}\simeq n_s\) has relative error \(1/N_s\).  Define the
shell's mass-weighted mean radial velocity and one-dimensional orbital
velocity dispersion by
\[
 M_s=\sum_{i\in s}m_i,\qquad
 \overline w_s=\frac{1}{M_s}\sum_{i\in s}m_iw_i,        \tag{PIC14a}
\]
\[
 \sigma_{v,s}^{\,2}=\frac{1}{3M_s}
 \sum_{i\in s}m_i\left[
   (w_i-\overline w_s)^2+\frac{2\eta_i}{r_i^2}
 \right].                                               \tag{PIC14b}
\]
The symbols \(r_i,w_i,\eta_i\) are marker radius, radial velocity, and half
the squared angular momentum.  Therefore \(2\eta_i/r_i^2\) is the squared
tangential speed.  The factor \(1/3\) converts the total three-dimensional
random kinetic variance into a one-dimensional dispersion.

For a sampled unordered pair \(i,j\), let
\(u_{ij}=|\bm v_i-\bm v_j|\) be relative speed and
\(\Sigma_{ij}(u_{ij})\) its physical pair area.  The tangential directions
are not stored in spherical phase space, so VENUS samples their relative
azimuth deterministically.  It then evaluates the collision-rate-weighted
mean free path directly from the evolving distribution:
\[
 \boxed{\displaystyle
 \lambda_s=
 \frac{\langle u_{ij}\rangle_s}
 {n_{{\rm partner},s}\left\langle
 u_{ij}\Sigma_{ij}(u_{ij})\right\rangle_s}}.             \tag{PIC14c}
\]
Angle brackets denote an average over unordered marker pairs and relative
tangential azimuth.  All pairs are used when the shell contains at most 512
unordered pairs; otherwise 512 deterministic pair samples are used.  For a material cross section per
mass \(s(u)=\sigma(u)/m\),
\[
 \Sigma_{ij}(u)=\frac{m_i+m_j}{2}s(u).
\]
For \(N_s\) equal masses and constant \(s\), equation (PIC14c) gives
\[
 \lambda_s=\frac{N_s}{N_s-1}\frac{1}{\rho_s s}.
\]
It approaches Frank and Tokayer's continuum result
\(\lambda_s=1/(\rho_s s)\) as \(N_s\) grows.  Retaining the finite-\(N_s\)
factor makes the locality diagnostic exactly consistent with the
unordered-pair trial law in equation (PIC17), including in a sparse shell.
For a velocity-dependent interaction, \(\sigma_{v,s}\) describes how broad
the evolving velocity distribution is, while equation (PIC14c) uses the full
sampled distribution rather than replacing it by one guessed RMS speed.  This
distinction matters when a low-speed tail has a much larger cross section.

If the shell width exceeds \(\lambda_s\), the shell is divided into two
equal-population children.  The criterion is applied recursively to each child.  Refinement
continues down to the configured minimum population.  The history output
records
\[
 \max_s \frac{r_{s,+}-r_{s,-}}{\lambda_s}.               \tag{PIC15}
\]
A value below one verifies that an accepted event joins markers within one
mean free path.  A production run can request a hard stop when equation
(PIC15) reaches one.

The source code implements this construction directly:
\codesource{WPIC adaptive collision-shell builder}
{\path{cpu/include/ev/weighted_vlasov_pic_collisions.hpp}}
{Measures the velocity-averaged mean free path and recursively divides a
 radial shell when its width is not local}
\begin{lstlisting}[language=C++]
const auto mfp = velocity_averaged_mean_free_path(
    markers, radius_order, begin, end, shell.volume,
    parameters, collision_options);
shell.mean_free_path = mfp.mean_free_path;
shell.orbital_velocity_dispersion =
    mfp.orbital_velocity_dispersion;

if (shell_width > shell.mean_free_path
    && shell_count >= 2 * minimum_shell_markers) {
    build_shell(left_half);
    build_shell(right_half);
}
\end{lstlisting}

\subsubsection{Step 6: choose the collision pairs}

Every pair of markers in one shell is physically eligible.  VENUS samples the
realized events with Bird's no-time-counter method.  Define
\[
 {\cal K}_{ij}(u)=u\Sigma_{ij}(u),
 \qquad {\cal K}_{\max}\geq
 \max_{a,b\in s}{\cal K}_{ab}(|\bm v_a-\bm v_b|).       \tag{PIC16}
\]
The symbol \({\cal K}_{ij}\) is the pair collision kernel: relative speed
times physical pair area.  It has dimensions of volume per time.
A Poisson process is a random event-counting model with independent waiting
intervals: for a constant rate \(\nu\), the chance of no event during
\(\Delta t\) is \(e^{-\nu\Delta t}\), and the waiting time to the next event
has an exponential distribution.  This statistical use of ``Poisson'' is
different from the Poisson field equation used for gravity.  The trial-event
clock is such a process with rate
\[
 \nu_{\rm trial}=\frac{N_s-1}{2}\,n_s{\cal K}_{\max}.    \tag{PIC17}
\]
If \(z_1\) is uniform on \((0,1)\), the time to the next trial is
\[
 \delta t_{\rm trial}=-\frac{\ln z_1}{\nu_{\rm trial}}. \tag{PIC18}
\]
For that trial, two distinct marker indices are drawn uniformly from the
shell.  With \(u=|\bm v_a-\bm v_b|\), a second uniform number \(z_2\) accepts
the pair when
\[
 z_2 < \frac{{\cal K}_{ab}(u)}{{\cal K}_{\max}}.         \tag{PIC19}
\]
Equations (PIC17)--(PIC19) make the accepted-pair rate proportional to
\(f(\bm v_a)f(\bm v_b)u\Sigma_{ab}(u)\).  The calculation visits the realized
trials and scales as \(N_p\log N_p+N_{\rm trial}\), including the radial sort.

The corresponding code is:
\codesource{WPIC no-time-counter collision sampler}
{\path{cpu/include/ev/weighted_vlasov_pic_collisions.hpp}}
{Draws trial times and distinct local marker pairs, then accepts each pair
 against the verified relative-speed--cross-section majorant}
\begin{lstlisting}[language=C++]
clock += -std::log(rng.uniform_open()) / trial_rate;
ia = rng.uniform_index(shell_count);
ib = distinct_uniform_index(shell_count, ia);

u_sigma = relative_speed * pair_total_cross_section(
    relative_speed, first_mass, second_mass, ...);
if (rng.uniform_open() * majorant_u_sigma < u_sigma) {
    scatter_pair(ia, ib);
}
\end{lstlisting}

Active markers may have unequal represented masses.  Each accepted elastic
event acts on the two complete masses in their mass-weighted centre-of-mass
frame.  Neither marker is split and both masses remain unchanged.

\subsubsection{Step 7: draw the outgoing direction}

A reduced spherical marker stores \(w=v_r\) and \(\eta=L^2/2\).  At its
current radius, its tangential speed is
\[
 v_t=\frac{\sqrt{2\eta}}{r}.                            \tag{PIC20}
\]
Spherical symmetry makes the tangential azimuth uniformly distributed.
VENUS draws an azimuth \(\alpha\) and lifts the marker to
\[
 \bm v=(w,v_t\cos\alpha,v_t\sin\alpha).                \tag{PIC21}
\]

At the exact coordinate origin, a finite spherical distribution has
\(\eta=0\). The stored \(|w|\) is the speed, and its Cartesian direction
is sampled isotropically. After scattering, the full outgoing speed is
stored with zero angular momentum, so the complete kinetic energy is
retained. The Cartesian pair update still conserves its vector momentum.
The mean projection along any reference axis is zero and its second moment
is \(w^2/3\). Shell dispersions, mean-free-path estimates, and collision
clocks use this same isotropic distribution. A positive radius whose square
underflows cannot represent a newly generated transverse energy in
\(\eta=r^2v_t^2/2\); the material operator rejects that state before any
collision is applied.
For an accepted pair, define
\[
 M_{12}=m_1+m_2,\qquad
 \bm V=\frac{m_1\bm v_1+m_2\bm v_2}{M_{12}},
 \qquad \bm u=\bm v_1-\bm v_2.                         \tag{PIC22}
\]
Here \(m_1,m_2\) are the complete marker masses, \(M_{12}\) is their sum,
\(\bm V\) is their centre-of-mass velocity, and \(u=|\bm u|\).

The differential cross section selects the direction of \(\bm u'\).  Write
\(\mu=\cos\theta\), where \(\theta\) is measured from \(\bm u\), and let
\(\psi\) be the azimuth around \(\bm u\).  For isotropic scattering, two
uniform numbers \(z_3,z_4\) give
\[
 \mu=2z_3-1,
 \qquad \psi=2\pi z_4.                                 \tag{PIC23}
\]

For the Yukawa--Born kernel, the numerical parameters \(m_\chi\) and
\(m_\phi\) set the speed scale \(m_\phi/m_\chi\), and
\(u=|\bm u|\) is the pair's relative speed defined in equation (PIC22).
Section~\ref{sec:yukawa-material-units} derives the amplitude and speed
conversion from physical particle masses and coupling to material-SIDM code
units.
With the abbreviations \(a=m_\phi^2\) and
\(b=m_\chi^2u^2/2\),
\[
 \frac{\dd\sigma}{\dd\Omega}\propto
 \frac{1}{[a+b(1-\mu)]^2},
 \quad a=m_\phi^2,
 \quad b=\frac{m_\chi^2u^2}{2}.                       \tag{PIC24}
\]
VENUS integrates this expression analytically.  A uniform \(z_3\) samples
the probability above the selected cosine.  With
\(y=(m_\chi u/m_\phi)^2\), the stable inverse is
\[
 \mu=1-\frac{2z_3}{1+y(1-z_3)}.                         \tag{PIC25}
\]
It approaches the uniform cosine \(1-2z_3\) as \(y\to0\).  At large
\(y\), the implementation divides numerator and denominator by \(y\),
keeping the arithmetic bounded.  The same inverse is used by material WPIC
and SphDSMC.
\codesource{Yukawa scattering angle}
{\path{cpu/include/ev/finite_angle_boltzmann.hpp}}
{Actual branches of \code{bp_fa_yukawa_born_mu_from_unit}; \code{z} is the
 sampled probability and \code{y} is the squared relative speed ratio}
\begin{lstlisting}[language=C++]
if (y <= 1.0) {
    one_minus_mu = 2.0 * z / (1.0 + y * (1.0 - z));
} else {
    const double inverse_y = 1.0 / y;
    one_minus_mu = 2.0 * z * inverse_y / (inverse_y + (1.0 - z));
}
\end{lstlisting}
Equation (PIC25) draws more small-angle events when the differential cross
section is forward peaked.  The azimuth remains \(\psi=2\pi z_4\).

An orthonormal basis \((\widehat{\bm u},\bm e_1,\bm e_2)\) then gives
\[
 \widehat{\bm u}'=
 \mu\widehat{\bm u}
 +\sqrt{1-\mu^2}\left(\cos\psi\,\bm e_1
 +\sin\psi\,\bm e_2\right).                           \tag{PIC26}
\]
The outgoing velocities are
\[
 \bm v_1'=\bm V+\frac{m_2}{M_{12}}u\widehat{\bm u}',
 \qquad
 \bm v_2'=\bm V-\frac{m_1}{M_{12}}u\widehat{\bm u}'.  \tag{PIC27}
\]
Thus each accepted event preserves all three components of pair momentum and
the pair kinetic energy.  The reduced marker variables are updated with
\[
 w_p'=v_{p,r}',
 \qquad
 \eta_p'=\frac{r_p^2}{2}\left[(v_{p,y}')^2+(v_{p,z}')^2\right]. \tag{PIC28}
\]
Equation (PIC28) is the angular momentum redistribution.  Each collision can
move both markers to new values of \(L=\sqrt{2\eta}\).

The main source lines are:
\codesource{WPIC finite-angle elastic scatter}
{\path{cpu/include/ev/weighted_vlasov_pic_collisions.hpp}}
{Samples the outgoing relative direction, applies the unequal-mass
 centre-of-mass map, and returns the updated spherical marker variables}
\begin{lstlisting}[language=C++]
const double total_mass = first_mass + second_mass;
const BPFAVec3 centre = bp_fa_scale(
    bp_fa_add(
        bp_fa_scale(v1, first_mass),
        bp_fa_scale(v2, second_mass)),
    1.0 / total_mass);
const double mu = sample_scatter_mu(u, p, options, rng);
const double psi = 2.0 * PI * rng.uniform_open();
const BPFAVec3 relative_after =
    bp_fa_scale(
        scattered_relative_direction(relative, mu, psi), u);

out1 = bp_fa_add(
    centre,
    bp_fa_scale(relative_after, second_mass / total_mass));
out2 = bp_fa_sub(
    centre,
    bp_fa_scale(relative_after, first_mass / total_mass));
marker.eta = 0.5 * r * r * (out.y * out.y + out.z * out.z);
\end{lstlisting}

\subsubsection{Step 8: schedule orbital motion, field refreshes, and collisions}

The frozen-field event library keeps separate gravitational field and
collision clocks.  The first clock measures the elapsed time since the field
was refreshed.  In \path{weighted_pic_boltzmann_potential_timescale}, its
deadline uses the represented-mass-weighted root-mean-square field and rate:
\[
 T_\Phi=\eta_\Phi\min\left[
 \left(\frac{\sum_i\omega_i\Phi_i^2}
             {\sum_i\omega_i\dot\Phi_i^2}\right)^{1/2},
 \left(\frac{\sum_i\omega_i g_i^2}
             {\sum_i\omega_i\dot g_i^2}\right)^{1/2}
 \right].                                             \tag{PIC29}
\]
Here \(M_i\geq0\) is deposited shell mass,
\(\omega_i=M_i/\sum_jM_j\), \(g=-\Phi_r\) is the radial acceleration defined
in (PIC12), and a dot denotes the physical-time derivative calculated from
the current moments by \code{vp_potential_time_derivative}.  The source is
\path{cpu/include/ev/weighted_pic_boltzmann_event.hpp}, lines 145--203.
Each channel contributes only when both of its mean squares exceed
\((4096\epsilon_{\rm mach})^2\); otherwise its time is infinite, where
\(\epsilon_{\rm mach}\) is double-precision machine roundoff.  Zero represented
mass and the helper's expanding-background branch also return infinity.
The standalone static/expanding KDK drivers instead use the successive-field
maximum-norm clock in (PIC39a)--(PIC39c).
The second clock measures the elapsed time since the last collision update.
Its deadline is
 \[
 T_{\rm coll}=\eta_{\rm coll}\min_r
 \left[n_{\rm partner}(r)\left\langle
 u\Sigma_{ij}(u)\right\rangle_r\right]^{-1}.          \tag{PIC30}
\]
Here \(\eta_\Phi\) and \(\eta_{\rm coll}\) are accuracy parameters no larger
than one, and \(n_{\rm partner}(r)=(N_s-1)/V_s\) is the finite-shell partner
density at the shell containing \(r\).  The averages in equation (PIC30) are
measured from the current markers in the local collision shells.  The default
pilot uses 128 uniformly chosen pairs per shell.  For each pair it draws the
unresolved tangential azimuths and evaluates \(u\Sigma_{ij}(u)\) with the
actual two masses.  If the sample mean is \(\overline{\cal K}_s\)
and its standard error is \(e_s\), the scheduling rate uses
\[
 {\cal K}_{s,\rm schedule}=
 \min({\cal K}_{s,\max},\overline{\cal K}_s+4e_s).      \tag{PIC30a}
\]
The four-error buffer makes an accidentally long collision interval unlikely.
The exact majorant \({\cal K}_{s,\max}\) from equation (PIC16) still controls every
internal NTC substep.  Thus the outer clock follows the physical mean rate,
while the event generator retains its strict upper bound.

Let \(A_\Phi\) and \(A_{\rm coll}\) be the two elapsed clock ages.  The next
frozen-field orbit interval is
\[
 \Delta t=\min\left(T_\Phi-A_\Phi,
                    T_{\rm coll}-A_{\rm coll},
                    \Delta t_{\rm orbit,max}\right).  \tag{PIC31}
\]
The orbit integrator may take many internal steps during \(\Delta t\).  These
internal steps only move the markers.  The endpoint distribution is
reconstructed after a scheduled event.  This arrangement permits accurate
pericentre integration, including Levi--Civita retries, while keeping the
number of phase space remaps small.

The finite cap \(\Delta t_{\rm orbit,max}\) is also a clock-probe interval.
If the cap arrives before either physical deadline, VENUS deposits only the
radial mass and velocity moments needed to remeasure \(T_\Phi\) and
\(T_{\rm coll}\).  It keeps the current gravitational field and does not
construct the full phase space array at this probe.  If a newly measured
deadline is already shorter than its accumulated age, the corresponding
event is carried out immediately.  This catches a rapid departure from an
initial equilibrium without introducing an extra remap at every orbit step.

At the end of the interval, VENUS carries out the event whose deadline has
arrived:
\begin{enumerate}
\item If the collision deadline arrives first, move the markers through the
frozen potential, apply the local NTC collisions from Steps 5--7, and then
reconstruct the endpoint distribution.
\item If the potential deadline arrives first, move the markers, reconstruct
the endpoint distribution, deposit its mass, and solve for the new potential.
The collision age is retained, so the pending collision deadline continues
to approach.
\item If both deadlines arrive together and \(T_{\rm coll}<T_\Phi\), use the
collision-first order in item 1.
\item If both deadlines arrive together and \(T_\Phi\leq T_{\rm coll}\),
reconstruct and refresh the field first, and then apply the collision update.
\end{enumerate}

Only the clock associated with a completed event is reset.  For example, a
field refresh leaves \(A_{\rm coll}\) unchanged.  Several field refreshes can
therefore occur during one collision interval, and several collision updates
can occur during one field-refresh interval.  After each reconstruction,
VENUS measures the current timescales again and plans the next event with the
remaining clock ages.

The integrated state keeps both continuous markers and a reconstructed
phase space snapshot.  The snapshot contains a cell-average grid and, when
needed, a short list of continuous carrier markers described in Step 9.  A
Boolean flag records whether this combined snapshot describes the current
markers.  A collision-first event ends with reconstruction, so the flag is
true.  At a coincident potential-first event, VENUS reconstructs the snapshot,
solves for the field, and then scatters the markers.  The local collision does
not change any marker radius or the radial mass profile, so the new field is
still valid.  The velocity space snapshot is marked stale and is rebuilt only
when phase space output or the next grid synchronization requires it.  This
uses one velocity reconstruction at that endpoint.

The source expresses this rule directly:
\codesource{WPIC two-clock event scheduler}
{\path{cpu/include/ev/weighted_vlasov_pic_event_scheduler.hpp}}
{Chooses the next orbit interval and orders coincident field-refresh and
 collision events without resetting the other physical clock}
\begin{lstlisting}[language=C++]
WeightedPICEventPlan event = plan_weighted_pic_event(
    clock, t_phi, t_collision,
    eta_phi, eta_collision, orbit_cap);

execute_weighted_pic_event(
    event,
    [&] { advance_markers_in_frozen_field(event.orbit_interval); },
    [&] { reconstruct_endpoint_distribution(); },
    [&] { apply_local_ntc_collisions(collision_age); });
commit_weighted_pic_event(clock, event);
\end{lstlisting}

The production wrapper calls the same operations through one function:
\codesource{WPIC Boltzmann event wrapper}
{\path{cpu/include/ev/weighted_pic_boltzmann_event.hpp}}
{Assembles transport, deposition, collision sampling, phase space
 reconstruction, and event diagnostics into one accepted update}
\begin{lstlisting}[language=C++]
WeightedPICBoltzmannEventDiagnostics diagnostic =
    advance_weighted_pic_boltzmann_event(
        state, field_grid, phase_geometry, parameters,
        orbit_plan, event_options, collision_options,
        remaining_run_time);
\end{lstlisting}
The returned diagnostic states the selected ordering, the number of NTC
trials and accepted pairs, whether reconstruction occurred before or after
collision, and the marker mass before and after the event.

The collision sampler state and both clock ages belong in a restart.  This
preserves the physical time remaining to each event after a continuation.

\subsubsection{\texorpdfstring{Step 9: reconstruct cell-average \(f(r,w,\eta)\)}
{Step 9: reconstruct cell-average f(r,w,eta)}}

The marker ensemble is itself a sampled distribution function.  Gravity needs
the radial deposition in equations (PIC7)--(PIC11).  A phase space snapshot is
reconstructed at a scheduled endpoint, an output, or a restart conversion.
The event order in Step 8 determines whether that endpoint reconstruction is
performed before or after the NTC collision update.

VENUS first places each active marker in one radial cell, then groups markers
with exactly the same stored mass.  ``Exactly'' here means literal equality
of the stored floating-point mass, not membership in a tolerance bin; the
remap never combines two unequal packet masses into one projection pair.
Within each numerical projection group it sorts by kinetic energy evaluated
at the cell center.  The lowest and highest remaining energies form a pair;
marker IDs break exact ties.  Pairing across the occupied energy range gives
the nonnegative projection a broad moment interval while each member still
uses velocity cells near its own continuous velocity.  This grouping is only
a reconstruction device: the physical NTC collision sampler remains free to
pair different masses.  For group mass \(m_g\), one continuous marker pair
has target radial momentum and kinetic energy
\[
 P_*=m_g(w_1+w_2),
 \qquad
 K_*=m_g\left[\frac{w_1^2+w_2^2}{2}
 +\frac{\eta_1+\eta_2}{r_i^2}\right].                 \tag{PIC32}
\]
The remap finds nonnegative weights \(a_\ell\) on nearby pairs of velocity
cells such that
\[
 \sum_\ell a_\ell=1,
 \quad
 \sum_\ell a_\ell(w_{\ell1}+w_{\ell2})=P_*/m_g,
 \quad
 \sum_\ell a_\ell(E_{\ell1}+E_{\ell2})=K_*/m_g.       \tag{PIC33}
\]
It deposits \(m_ga_\ell\) into each member of pair \(\ell\).  An odd marker,
including a marker whose mass is unique in that radial cell, is paired with
an identical copy and each copy receives half the ordinary pair weight.  This
does not split or add a persistent marker; it is only a two-point
representation used while depositing one unchanged mass \(m_i\).  Thus a
discrete unequal-weight hierarchy can form ordinary equal-mass projection
groups, whereas a realistic continuously varying mass spectrum may leave
almost every group as a singleton.  The latter case is algebraically valid
but can be expensive and noisy; record the duplicated-singleton fraction,
projection cost, and carrier fraction, and refine both the velocity grid and
marker population before interpreting it.  The reconstructed
cell average is
\[
 f_{ijk}=\frac{M_{ijk}}{
 8\pi^2\Delta r_i\Delta w_j\Delta\eta_k}.             \tag{PIC34}
\]
Equations (PIC33)--(PIC34) preserve active marker mass, radial momentum, and
kinetic energy while keeping every cell mass nonnegative.  Hidden-inner and
escaped marker masses remain separate ledger entries.  A marker pair outside
the velocity domain stops the remap and reports whether the \(w\) or \(L\)
extent must be enlarged.

A continuous velocity can also lie inside the grid faces while its exact
moments lie outside the set formed by the nearby cell centers.  For example,
consider a marker with \(w=0\) and \(\eta=0\) on a cell-average grid whose
first \(\eta\) center is positive.  The marker belongs to the grid domain, but
every nonnegative combination of cell-center packets has positive kinetic
energy.  Placing the marker in the first cell would therefore add kinetic
energy.

VENUS searches a bounded neighborhood of velocity cells for the nonnegative
weights in equation (PIC33).  The option
\code{maximum_candidate_radius} limits this neighborhood, so the cost of one
difficult pair stays local.  If an in-domain pair still has no solution, the
pair can remain in continuous form as a reconstruction carrier.  Its marker
indices are stored in \code{carrier_marker_indices}.  The reconstructed state
then obeys
\[
 M_{\rm active}=\sum_{ijk}M_{ijk}+M_{\rm carrier},       \tag{PIC34a}
\]
and the two other budgets obey
\[
 \begin{aligned}
 P_{r,\rm active}&=P_{r,\rm grid}+P_{r,\rm carrier},\\
 K_{\rm active}&=K_{\rm grid}+K_{\rm carrier}.
 \end{aligned}                                          \tag{PIC34b}
\]
The carrier markers remain in the main marker ensemble.  Orbit motion,
radial gravity deposition, and NTC pair selection continue to include them.
The carrier list records the small part of the endpoint distribution that the
current velocity-cell centers cannot represent with exact moments.

The option \code{maximum_carrier_mass_fraction} is a hard cumulative upper
bound on \(M_{\rm carrier}/M_{\rm active}\).  Before retaining any unresolved
pair, VENUS adds that pair's full represented mass to the carrier mass already
accepted in the same remap and checks the resulting fraction.  The remap stops
rather than silently crossing the cap.  Its default value is zero, so the
first unresolved projection is reported immediately.  A diagnostic run may
set a small positive bound and measure how this fraction changes as marker
count and velocity resolution increase.  A production convergence sequence
requires the fraction to decrease under refinement.  A velocity outside the
grid faces always stops the remap, regardless of the carrier bound.

The bridge is called with:
\codesource{WPIC conservative phase space remap}
{\path{cpu/include/ev/weighted_vlasov_pic_phase_remap.hpp}}
{Projects continuous weighted markers onto cell averages while reporting any
 unresolved carrier mass and moment-closure residuals}
\begin{lstlisting}[language=C++]
WeightedPICPhaseRemapOptions remap_options;
remap_options.maximum_candidate_radius = 8;
remap_options.maximum_carrier_mass_fraction = 0.01;

WeightedPICPhaseRemapResult remap =
    remap_weighted_vlasov_pic_to_phase_space(
        markers, geometry, parameters, remap_options);

double represented_mass = std::accumulate(
    remap.cell_mass.begin(), remap.cell_mass.end(), 0.0);
represented_mass += remap.diagnostics.carrier_mass;
\end{lstlisting}

The useful remap cadence follows the event intervals in equations
(PIC29)--(PIC31), which are usually longer than an orbit-integrator substep.
Characteristic motion accumulates continuously on the markers, and equation
(PIC34) is evaluated once at a scheduled endpoint.  Convergence studies vary
the marker count, both clock tolerances, \(r_{\min}\), and all three grid
resolutions.
An out-of-domain velocity is a hard domain error rather than an implicit
boundary remap.  Equations (PIC34a)--(PIC34b) provide the mass, momentum, and
kinetic-energy closure conditions for every reconstruction.

\subsubsection{Initial conditions, provenance, and task interfaces}

There are two fresh-data routes in this module, and they produce different
objects.

\noindent\textbf{Expanding Vlasov--Poisson and subhalo route.}
The consuming task is \code{venus_cpu_subhalo_wpic}.  Its fail-closed frontend
is \path{cpu/src/apps/subhalo/expanding_wpic.cpp}.  In the shared driver,
\code{PowerLawInitialCondition} is the source class that evaluates the
normalized, tapered power-law density selected by the cusp exponent
\(\gamma\), and precomputes the shell and enclosed masses from which markers
are sampled.  This class and the \code{initialize_markers} sampler in
\path{cpu/src/drivers/expanding_wpic.cpp} create an in-memory
\code{std::vector<WeightedPICMarker>}, meaning a resizable list of the marker
records defined at the start of this chapter.  The first accepted state is
persisted as \path{restart_state.bin} in the requested output directory.
Before using it, verify the supported cusp exponent and cosmological
convention, positive marker weights, represented mass, marker--mesh support,
the reported initial density error, and an otherwise identical zero-collision
control.  Then refine the marker loading and radial mesh.

\noindent\textbf{Static material-SIDM route.}
The consuming task is \code{venus_cpu_sidm_wpic_constant}, selected by
\path{cpu/src/apps/sidm/wpic_constant.cpp} and implemented by
\path{cpu/src/drivers/sidm_wpic.cpp}.  A compatible Frank marker population can
be generated by
\path{cpu/tools/generate_frank_core_spherical_dsmc_restart.py}; its output is a
versioned \code{VENUSSPHDSMC1} persistent-marker restart.  The importer
\code{import_weighted_pic_frank_restart} in
\path{cpu/include/ev/weighted_pic_frank_restart_import.hpp} converts
\((r,v_r,L,m,I)\) to WPIC markers \((r,w,\eta,m,I)\), after which the runner
writes its own \path{restart_state.bin}.  Validate the restart magic and
version, unique identifiers, positive individual masses, represented mass,
radial support, initial equilibrium drift, shell locality, and collision rate.
An unequal-mass input must be imported in full; a marker-limited import is only
an equal-mass diagnostic and is not a convergence datum.

\noindent\textbf{Runnable expanding WPIC example.}
The WPIC task executables use command-line controls rather than the generic
\path{params.ini}.  This self-contained example enables both elastic subhalo
scattering and inelastic capture:
\codesource{\code{venus_cpu_subhalo_wpic}}
{\path{cpu/src/apps/subhalo/expanding_wpic.cpp};
 \path{cpu/src/drivers/expanding_wpic.cpp}}
{Selects the expanding subhalo task and owns its command-line setup,
 marker initialization, evolution, and output}
\begin{lstlisting}[basicstyle=\ttfamily\scriptsize]
cd VENUS/cpu
make venus_cpu_subhalo_wpic
make venus_cpu_vlasov_poisson_wpic_expanding

./venus_cpu_subhalo_wpic \
  --gamma 0.5 \
  --end-time 0.2 \
  --maximum-timestep 0.01 \
  --particles-per-radial-cell 512 \
  --subhalo-collision-strength 0.1 \
  --subhalo-velocity-scale 0.1 \
  --subhalo-plummer-area-per-mass 1 \
  --subhalo-low-speed-regularization 0.05 \
  --subhalo-collision-interval 0.01 \
  --subhalo-collision-timescale-fraction 0.1 \
  --subhalo-collision-cfl 0.2 \
  --subhalo-collision-maximum-substeps 256 \
  --threads 8 \
  --output-directory output/wpic_subhalo_example
\end{lstlisting}
For every strength, cusp, mesh, and marker loading, inspect the evolving
locality diagnostic before interpreting the profiles.  Set
\code{--subhalo-collision-strength 0} only in the compatibility driver for an
otherwise identical collisionless Vlasov--Poisson control; the task-specific
\code{venus_cpu_subhalo_wpic} frontend requires positive collision strength.
The static elastic SIDM runner can start from a generated equal-mass Frank
marker sample, while its importer and collision core also preserve unequal
masses stored by version-4/5 SphDSMC restarts; the runnable import is in
Section~\ref{sec:wpic-material-sidm-run}.  The value 256 in this short example
is only the administrative internal-loop guard.  High-rate unequal-mass
production uses 16384; the collision CFL still selects every physical NTC
substep.

\subsubsection{Running the integrated material-SIDM core case}
\label{sec:wpic-material-sidm-run}

The task-specific executable is
\code{venus_cpu_sidm_wpic_constant}.  It keeps the
markers, field, two clocks, collision sampler, and reconstructed snapshot in
one integrated event state.  Every step calls the high-level event advance
function shown in Step 8, so transport, collision sampling, reconstruction,
and conservation accounting share one event order.

The following commands generate a compatible equal-mass Frank marker object,
import all of its markers, and advance a short constant-cross-section
interval.  A gridded \code{EVLOGRST1} file is not accepted by
\code{--frank-restart}; that option requires the persistent-marker
\code{VENUSSPHDSMC1} format.
\codesource{\path{generate_frank_core_spherical_dsmc_restart.py} and
 \code{venus_cpu_sidm_wpic_constant}}
{\path{cpu/tools/generate_frank_core_spherical_dsmc_restart.py};
 \path{cpu/src/apps/sidm/wpic_constant.cpp};
 \path{cpu/src/drivers/sidm_wpic.cpp}}
{Generates the versioned Frank marker population, imports it into the
 constant-cross-section material-SIDM WPIC task, and advances the coupled
 field--collision state}
\begin{lstlisting}
python tools/generate_frank_core_spherical_dsmc_restart.py \
  --output output/frank_sphdsmc_50000.restart.bin \
  --particles 50000 --seed 5489 --radial-sampling stratified

make venus_cpu_sidm_wpic_constant
./venus_cpu_sidm_wpic_constant \
  --frank-restart output/frank_sphdsmc_50000.restart.bin \
  --output-directory output/sidm_wpic_constant \
  --end-time 0.0033 \
  --maximum-timestep 0.001 \
  --potential-event-eta 1.0 \
  --collision-event-eta 0.01 \
  --radial-cells 64 \
  --radial-velocity-cells 32 \
  --eta-cells 96 \
  --eta-grid-power 8 \
  --radial-velocity-maximum 5 \
  --angular momentum-maximum 400 \
  --sidm-cross-section constant \
  --sidm-sigma0 20.106192982974676 \
  --collision-strength 1 \
  --collision-cfl 0.2 \
  --collision-maximum-substeps 65536 \
  --collision-adaptive-shell-timestep \
  --collision-adaptive-maximum-interval 0.02 \
  --remap-candidate-radius 8 \
  --maximum-carrier-mass-fraction 0.01 \
  --threads 32
\end{lstlisting}
The optional \code{--marker-limit} performs a deterministic radial subsample
of the supplied equal-mass realization, then rescales it to
the full input mass.  This option is intended only for inexpensive inspection;
omitting it imports every marker with its
stored mass unchanged.  A limited import of an unequal-mass restart is
rejected because number-ranked subsampling could bias its phase space mass
profile; import all such markers.  The limited run's core-density diagnostic
is not quantitatively comparable with the full import or with a different
marker limit.  The core is defined by the innermost
\(\max[250,\operatorname{round}(2.5\times10^{-4}N_{\rm selected})]\)
selected markers.  Changing \(N_{\rm selected}\), while rescaling each
selected marker to represent the full mass, changes the enclosed mass and
radius that define ``core.''  Consequently \code{core_density},
\code{core_sigma_r2}, and \code{core_radius} from a marker-limited history
cannot be used as a physical convergence comparison.

The radial profiles use the conservative deposition in equations
(PIC7)--(PIC11) directly.  Every scheduled phase space reconstruction is
recorded in
\path{operator_mass_budget.csv}.  Its columns include the event order, marker
mass before and after, gridded mass, carrier mass, radial-momentum error,
kinetic-energy error, and maximum pair-projection error.  A final output
synchronization has event order \(-1\).

The restart stores the continuous markers, collision random-number state,
both event-clock ages, the frozen Poisson field, and the two current timescale
estimates. A continuation starts with \code{--pic-restart}. Restoring the
saved field preserves the potential interval already in progress. The next
collision or potential event then occurs at the same time as in uninterrupted
evolution.
\codesource{Static SIDM event continuation}
{\path{cpu/src/drivers/sidm_wpic.cpp};
 \path{cpu/include/ev/weighted_pic_boltzmann_event.hpp}}
{Retain the frozen field and the current physical event deadlines}
\begin{lstlisting}[language=C++]
state.field = checkpoint.frozen_field;
state.potential_timescale = checkpoint.potential_timescale_estimate;
state.collision_timescale = checkpoint.collision_timescale_estimate;
\end{lstlisting}
The saved control record authenticates the mesh, force and collision laws,
orbit tolerances, event controls, and thread count. End time and output
cadence may change on continuation. For a deliberate mesh or numerical-control change,
\code{--refresh-restart-field} recomputes the field and its timescales and
starts a new potential interval. It retains the collision age accumulated
since the last collision, so the physical gravity and collision laws retain
their saved values. The driver records this choice in
\path{restart_import.txt}. Legacy versions 1--6 contain particle initial data;
their import reconstructs the field and timescales and reports that operation.

Continuing in an existing output directory checks every saved history and
profile time against the selected checkpoint before opening output files.
If a saved diagnostic is later than that checkpoint, use a separate output
directory for the continuation. This preserves the original trajectory and
keeps each recorded history in time order. The expanding WPIC driver applies
the same check to its energy, virial, and collision histories.

\paragraph{Particle export and numerical-sample refinement.}
The inspection tool has an explicit conversion command,
\par\noindent\begin{minipage}{\linewidth}
\codesource{Export particle initial data}
{\path{cpu/tools/inspect_weighted_pic_restart.cpp}}
{Preserve the particle state and initialize a fresh field schedule}
\begin{lstlisting}[language=bash]
inspect_weighted_pic_restart INPUT --export-particles OUTPUT
\end{lstlisting}
\end{minipage}\par
It preserves particle masses, coordinates, internal state, identifiers,
random counters, time, and step. Its version-6 output initializes a new field
and event schedule at that physical time, with zero event ages and fresh
timescale estimates. This supplies particle initial data for a receiving
CPU or GPU task. Export requires a completed collision interval,
\code{collision_event_age=0}, and a distinct output file.

The refinement command divides each numerical marker into equal-mass
samples at the same phase space point:
\codesource{Refine the numerical particle sampling}
{\path{cpu/tools/refine_weighted_pic_restart.cpp}}
{Split numerical weights while preserving the represented distribution}
\begin{lstlisting}[language=bash]
refine_weighted_pic_restart INPUT OUTPUT FACTOR --reinitialize-field
\end{lstlisting}
The factor is an integer from 2 through 16. Refinement preserves the
mass-weighted spatial and velocity moments and assigns distinct identifiers
to the additional samples. It requires zero internal energy, zero internal
velocity dispersion, and a completed collision interval. These conditions
support collisionless or material-SIDM sampling studies. A subhalo
marker represents a physical object whose mass enters its encounter law;
such an object's mass distribution must be specified physically. A
frozen-field input requires the explicit \code{--reinitialize-field} option.
The output is version-6 particle initialization with the frozen field and
timescale estimates cleared and the potential-event age reset.

\paragraph{CPU material-SIDM local shell clocks.}
Let
\[
 \nu_s=n_{{\rm partner},s}
       \max_{i,j\in s}[u\Sigma_{ij}(u)]
\]
be the conservative NTC exposure rate in collision shell (s).  Reference
mode uses the fastest shell for every shell,
\[
 \delta t_{\rm fixed}=\min\left(
 \Delta t_{\rm op}-t,
 \frac{C_{\rm coll}}{\max_s\nu_s}\right).             \tag{PIC34c}
\]
With \code{--collision-adaptive-shell-timestep}, every shell has an
independent elapsed time \(t_s\) and instead takes
\[
 \delta t_s=\min\left(
 \Delta t_{\rm op}-t_s,
 \frac{C_{\rm coll}}{\nu_s}\right).                    \tag{PIC34d}
\]
The shells are disjoint and marker radii do not move inside the split
collision operator.  Threads therefore advance different shells concurrently
without changing their event rates.  After each local elastic update VENUS
refreshes the shell velocity moments and safe \(u\Sigma\) majorant; shell
faces, volume, density, and radial membership are reused.  The implementation
is structurally
\codesource{WPIC adaptive local shell clocks}
{\path{cpu/include/ev/weighted_vlasov_pic_collisions.hpp};
 \path{cpu/src/drivers/sidm_wpic.cpp}}
{Advances disjoint constant-cross-section shells with their own conservative
 collision exposure while retaining one synchronized outer event}
\begin{lstlisting}[language=C++]
for (shell : collision_shells) {
  while (shell.elapsed < operator_interval) {
    rate = shell.partner_density * shell.majorant_u_sigma;
    dt_shell = min(operator_interval - shell.elapsed,
                   collision_cfl / rate);
    apply_shell_ntc(shell, dt_shell, shell.local_substep);
    refresh_shell_velocity_majorant(shell);
  }
}
\end{lstlisting}
\code{--collision-adaptive-maximum-interval} caps the synchronized outer
collision batch; it is not a second local CFL.  In this CPU release the
reusable-shell path is deliberately restricted to constant cross section,
equal marker masses within each shell, and \(0<C_{\rm coll}\le1\).  The
velocity-dependent and unequal-mass material paths retain equation (PIC34c).
The reported \code{substeps} is the largest local count and
\code{shell_visits} is the sum over shells.  Omitting the flag retains the
mesh-wide reference schedule.

\subsubsection{Static and expanding self-consistent KDK}
\label{sec:wpic-expanding-kdk}

The expanding runner stores physical radius \(r=a(t)x\) and physical peculiar
radial velocity \(u\), where \(a(t)\) is the cosmological scale factor and
\(x\) is comoving radius.  The canonical radial momentum is
\[
 p=a u.
\]
For a power-law background
\[
 a(t)=a_0\left(\frac{\tau}{\tau_0}\right)^\alpha,
 \qquad \tau=t+t_{\rm off},
\]
\(\alpha\) is the expansion exponent and \(t_{\rm off}\) is the configured
time offset.  During a force-free drift the useful interval is
\[
 s(t_0,t_1)=\int_{t_0}^{t_1}\frac{\dd t}{a(t)^2}.
\]
Let \(\tau_0=t_0+t_{\rm off}\), \(\tau_1=t_1+t_{\rm off}\), and
\(\beta=1-2\alpha\).  The exact integral is
\[
 s=\frac{\tau_0}{a(t_0)^2}
 \begin{cases}
  \displaystyle
  \frac{\exp\!\left[\beta\ln(\tau_1/\tau_0)\right]-1}{\beta},
  &\beta\ne0,\\[8pt]
  \ln(\tau_1/\tau_0),&\beta=0.
 \end{cases}
\]
The code evaluates the numerator with \code{expm1} to avoid cancellation when
\(|\beta\ln(\tau_1/\tau_0)|\ll1\).

At fixed specific angular momentum \(L\), the exact free central drift in
canonical comoving variables is
\[
 x_1^2=(x_0+s p_0)^2+
       \left(\frac{sL}{x_0}\right)^2,
\]
\[
 p_1=\frac{p_0(x_0+s p_0)
 +(L/x_0)(sL/x_0)}{x_1}.
\]
These equations are the radial projection of a straight Cartesian comoving
trajectory and therefore resolve the centrifugal turning point without
integrating the singular-looking term \(L^2/r^3\) explicitly.  The physical
endpoint variables are
\[
 r_1=a(t_1)x_1,\qquad u_1=\frac{p_1}{a(t_1)}.
\]

With peculiar potential \(\Phi\), the canonical force is
\[
 F_p=-a\Phi_r.
\]
The self-consistent KDK sequence is
\[
 p_{1/2}=a_0u_0+\frac{\Delta t}{2}F_p(t_0),
\]
followed by the exact drift above, endpoint deposition and Poisson solve, then
\[
 p_1\leftarrow p_1+\frac{\Delta t}{2}F_p(t_1),
 \qquad u_1=p_1/a_1.
\]
In a static universe \(a=1\) and \(s=\Delta t\), so the same formulas reduce
to an exact free central drift between ordinary gravitational half kicks.

The grid boundaries remain at fixed physical radii. A reflecting wall at
\(R\) reverses the physical radial velocity at its crossing time
\(t_{\rm hit}\). In this subsection's peculiar-velocity notation,
\(u_{\rm out}=-u_{\rm in}-2H(t_{\rm hit})R\). The tangential velocity
\(L/R\) remains unchanged, so the collision conserves kinetic energy at
the wall. The orbit then follows the free drift for the remaining interval.

To locate every crossing, introduce the dimensionless drift coordinate
\(h=a(t_0)^2s/\tau_0\). The radius during the free drift is
\[
 r(h)=\frac{a(t)}{a(t_0)}
 \left[(r_0+\tau_0u_0h)^2+
       \left(\frac{\tau_0Lh}{r_0}\right)^2\right]^{1/2}. \tag{PIC34h}
\]
The relation between \(h\) and time follows from the integral for \(s\)
above. Differentiating (PIC34h) shows that its radial turning points satisfy
\[
 \begin{split}
 0={}&(1-\alpha)\tau_0^2
       \left[u_0^2+(L/r_0)^2\right]h^2\\
 &+\left[\tau_0r_0u_0+
        \tau_0^2\left[u_0^2+(L/r_0)^2\right]\right]h
   +\alpha r_0^2+\tau_0r_0u_0 .
 \end{split}
 \tag{PIC34i}
\]
The at most two roots divide the drift into intervals over which radius
changes in one direction. Bisection on each interval finds the first
intersection with either wall. This also detects an inward crossing followed
by an outward crossing during one step. After reflection the search repeats
over the remaining time; exceeding 64 impacts raises an error.
\codesource{Specular reflection during the free angular drift}
{\path{cpu/include/ev/central_drift.hpp}}
{Locate each impact, reverse its physical normal velocity, and continue the unused interval}
\begin{lstlisting}[language=C++]
const State impact = at_h(r, w, L, tau, alpha, hit_h);
const double elapsed = tau * std::expm1(log_time_ratio(hit_h, alpha));
tau += elapsed;
r = wall;
w = -impact.w - 2.0 * alpha * wall / tau;
\end{lstlisting}
The shared CPU/CUDA helper accepts signed intervals, so the same collision
rule applies when tracing a characteristic backward. Static tests conserve
kinetic energy and recover the initial orbit under time reversal. Expanding
tests agree with an independent integration of the physical radial equations.
At the true origin, radial motion continues through the parity map, while
nonzero angular momentum is retained by the free central drift.


The public implementation is deliberately split:
\codesource{WPIC self-consistent kick--drift--kick step}
{\path{cpu/include/ev/weighted_vlasov_pic.hpp}}
{Uses the old field for the first half kick, solves gravity from the endpoint
 deposition, and uses that new field for the second half kick}
\begin{lstlisting}[language=C++]
begin_weighted_pic_self_consistent_kdk_parallel(
    markers, time, dt, q_grid, r_grid, field_n, p, threads);

endpoint = deposit_weighted_vlasov_pic_parallel(...);
field_np1 = solve_weighted_pic_poisson_field(...);

finish_weighted_pic_self_consistent_kdk_parallel(
    markers, time + dt, dt, q_grid, r_grid,
    field_np1, p, threads);
\end{lstlisting}
This makes the endpoint field solve visible and prevents an accidental use of
the old field in the second kick.  The current path requires KDK, active
markers, a reflecting outer boundary, and reflecting or origin-parity central
handling.  The generic delayed-return callback is available in the WPIC
library, but the public self-consistent adapter does not invent a first
crossing time for a delayed marker.

\subsubsection{Elastic material and gravitational kernels in WPIC}

The collision module accepts unequal positive marker masses and has two
physics selections:
\begin{itemize}
\item \code{venus_cross_section}: the constant, power-law, or Born-Yukawa
material cross sections also used by the VENUS Boltzmann and SphDSMC modules;
\item \code{gravitational_plummer}: a physical two-body deflection in a
Plummer-softened gravitational potential.
\end{itemize}
For the first selection, the configured law is
\[
 s(u)=\frac{\sigma(u)}{m},
\]
and the physical pair cross section and kernel are
\[
 \Sigma_{ij}(u)=\frac{m_i+m_j}{2}s(u),\qquad
 K_{ij}(u)=u\Sigma_{ij}(u).                               \tag{PIC26a}
\]
Here \(u\) is relative speed and \(m_i,m_j\) are the two complete partner
masses.  Constant and power-law kernels use isotropic
outgoing directions.  The Born-Yukawa kernel uses the forward-peaked inverse
cumulative distribution in equations (PIC24)--(PIC25).  The legacy scalar
\code{softened_gravity} transfer-cross-section fit is not accepted as a
differential scattering law in this module; a gravitational encounter must
use the explicit Plummer branch.

Shells are radius sorted with marker identifier as a stable tie breaker.
Their mass density is \(\sum_i m_i/V_s\), their number density is \(N_s/V_s\),
and their locality diagnostic uses the evolving pair-rate average in equation
(PIC14c), including both unequal partner masses and the actual
velocity-dependent cross section.  The larger NTC pair-kernel
\emph{majorant} is only a rejection-sampling bound: it controls how many trial
pairs are proposed and must not be substituted for a physical mean free path.
Distinct shells are processed independently, and thread results are reduced
in deterministic shell order.

\subsubsection{Elastic gravitational scattering by a softened Plummer pair}
\label{sec:wpic-plummer-scattering}

Let \(m_i,m_j\) be the two body masses, \(\epsilon_{ij}\) the pair softening
length, \(R\) their separation, \(b\) the
impact parameter at infinity, and \(u\) the relative speed at infinity.  The
relative potential per unit reduced mass is
\[
 \phi_{ij}(R)=-\frac{G(m_i+m_j)}
 {\sqrt{R^2+\epsilon_{ij}^2}}.
\]
The relative acceleration is
\[
 \ddot{\bm R}
 =-\frac{G(m_i+m_j)\bm R}
 {(R^2+\epsilon_{ij}^2)^{3/2}}.
\]
Conservation of relative energy and angular momentum gives the radial
function
\[
 F(R)=1-\frac{b^2}{R^2}
 +\frac{2G(m_i+m_j)}
 {u^2\sqrt{R^2+\epsilon_{ij}^2}}.
\]
The positive turning radius \(R_0\) satisfies \(F(R_0)=0\).  VENUS brackets
this root and finds it by bisection.  The deflection angle is
\[
 \Theta(b,u)
 =\left|2b\int_{R_0}^{\infty}
 \frac{\dd R}{R^2\sqrt{F(R)}}-\pi\right|.
\]
The change of variable \(R_0/R=1-z^2\) removes the square-root endpoint
singularity.  Adaptive Simpson quadrature in extended precision evaluates the
regularized integral to the configured deflection tolerance.  Recursion is
bounded by \path{gravitational_deflection_max_depth}; the corresponding
accuracy key in the collision options is
\path{gravitational_deflection_tolerance}.

Impact parameters are uniform in incident area:
\[
 b=\sqrt{b_{\min}^2+z(b_{\max}^2-b_{\min}^2)},
 \qquad 0<z<1.
\]
The configured \((b_{\min},b_{\max},\epsilon)\) refer to a body of reference
mass \(m_{\rm ref}\).  VENUS scales all encounter lengths by
\[
 \ell_{ij}=\sqrt{\frac{m_i+m_j}{2m_{\rm ref}}},
\]
so \(b_{ij}=\ell_{ij}b\), \(\epsilon_{ij}=\ell_{ij}\epsilon\), and the
physical pair aperture is
\[
 \Sigma_{{\rm area},ij}
 =A_{\rm grav}\frac{m_i+m_j}{2m_{\rm ref}}\,
 \pi(b_{\max}^2-b_{\min}^2).
\]
where \(A_{\rm grav}\) is the configured strength.  A finite \(b_{\max}\) is
required: gravity has a long-range weak-deflection tail, so the encounter
operator is undefined until the resolved impact-parameter interval is stated.
The sampled azimuth fixes the scattering plane and \(\Theta(b,u)\) rotates the
relative velocity inside that plane.

Two useful limits characterize the implementation.  As \(\epsilon/b\to0\), the
Rutherford point-mass angle is recovered.  In the weak-impulse limit,
\[
 \Theta\simeq
 \frac{2G(m_i+m_j)b}
 {u^2(b^2+\epsilon_{ij}^2)}.
\]
The exact quadrature is retained because a softened core can have a
non-monotone, or ``rainbow'', deflection function that is not represented by
the weak-angle formula.

\subsubsection{Resolved Plummer subhalo outcomes: scattering versus capture}
\label{sec:wpic-subhalo-capture}

The unequal-mass subhalo operator represents two mutually exclusive
outcomes of one resolved Plummer encounter:
\[
 {\cal C}={\cal C}_{\rm scatt}+{\cal C}_{\rm cap}.
\]
Here \({\cal C}_{\rm scatt}\) is elastic unbound scattering and
\({\cal C}_{\rm cap}\) is inelastic capture.  Let the two complete physical
subhalos have masses \(M_i,M_j\), Plummer radii \(a_i,a_j\), internal
velocity scales \(s_i,s_j\), relative velocity
\(\bm u=\bm v_i-\bm v_j\), and mass-weighted centre velocity
\[
 \bm V=\frac{M_i\bm v_i+M_j\bm v_j}{M_i+M_j}.
\]

For a resolved unequal-mass encounter with pericentre \(r_p\), compute the
internal response of the two objects separately and set
\[
 M_{\rm t}=M_i+M_j,\qquad
 \mu_{ij}=\frac{M_iM_j}{M_{\rm t}},\qquad
 \Delta E_{\rm int}=\Delta E_i+\Delta E_j.
\]
Capture occurs when this internal heating exceeds the
centre-of-mass-frame relative kinetic energy at infinity:
\[
 \Delta E_{\rm int}(u,r_p)\ge\frac12\mu_{ij}u^2.
\]
For \(M_i=M_j=M\), the right-hand side reduces to \(Mu^2/4\).
If \(b_\infty\) is the incident impact parameter, the focused capture area per
logarithmic pericentre interval is
\[
 \frac{\dd\Sigma_{\rm cap}}{\dd\ln r_p}
 =\pi r_p\frac{\dd b_\infty^2}{\dd r_p}
 H\!\left(\Delta E_{\rm int}-\frac12\mu_{ij}u^2\right),
\]
where \(H\) is the Heaviside step function, equal to one when its argument is
positive and zero when negative.  Gravitational focusing supplies
\[
 \frac{\dd b_\infty^2}{\dd r_p}
 =2r_p\left[
 1-\frac{2\phi_{12}(r_p)+r_p\phi'_{12}(r_p)}{u^2}
 \right].
\]
Here \(\phi_{12}\) is the relative two-subhalo potential per unit reduced mass
and a prime denotes differentiation with respect to \(r_p\).

The complementary unbound region contributes a momentum-transfer rather than
raw geometric cross section:
\[
 \Sigma_{{\rm T},{\rm unb}}
 =\int_0^{a_i+a_j}(1-\cos\Theta)
 \frac{\dd\Sigma_{\rm unb}}{\dd\ln r_p}\,\dd\ln r_p .
\]
The factor \(1-\cos\Theta\) weights a weak forward deflection less than a large
momentum-changing deflection.  Capture and scattering are therefore an
exclusive partition of the same focused calculation; they are not two
independently normalized fitting functions.

The table is calibrated by following a stationary harmonic-core ensemble in
the full spatial Plummer force of the other object; the two core responses
are evaluated separately.  Each independently sampled orbit geometry is
represented at four harmonic phases.  In addition, the exact first variation
of the discrete orbit map is propagated along the unperturbed trajectory.
Its energy contribution is reversible work, linear in the perturber force,
whose continuum phase average is zero.  Subtracting this finite-sample term
prevents it from overwhelming the physical quadratic heating at very small
mass ratio while retaining the nonlinear response.  The encounter table
stores the raw energy, the subtracted linear term, and
\(\max(\Delta E_{\rm raw}-\Delta E_{\rm linear},0)\) separately.  The maximum
with zero means that residual numerical cooling is treated as an unresolved
zero, not as negative irreversible heating.  The encounter-table loader
checks this identity in every speed--pericentre cell.  Nested particle-number
refinement is required so that this non-negative projection cannot create a
spurious capture channel.  The focused-orbit and first-variation construction
above is the derivation used by VENUS; the table definitions, identities,
and interpolation rules below form its complete release interface.

The equal-mass curve is only the \(q=1\) boundary of the physical kernel.
For two homologous Plummer objects define
\[
 q=\frac{\min(M_i,M_j)}{\max(M_i,M_j)},\qquad
 a_*=\frac{a_i+a_j}{2},\qquad
 x=\frac{u}{s_{ij}}.
\]
The calibrated source table contains 72 common speed nodes at every mass-ratio
node and stores
\[
 S_{\rm cap}(x;q)=\frac{\Sigma_{\rm cap}}{\pi a_*^2},
 \qquad
 S_{\rm scatt}(x;q)=
 \frac{\Sigma_{{\rm T},{\rm unb}}}{\pi a_*^2}.
\]
The accepted mass-ratio grid is
\[
 q=(10^{-6},10^{-5},10^{-4},10^{-3},3\times10^{-3},
 10^{-2},3\times10^{-2},0.1,0.2,0.5,1).
\]
The final nonzero capture node is \(x=0.796\) at equal mass,
\(0.502\) at \(q=10^{-2}\), \(0.240\) at \(q=10^{-3}\), and
\(0.087\) at \(q=10^{-4}\).  No capture occurs at any stored speed
\(x\ge0.05\) for \(q=10^{-5}\) or \(10^{-6}\).  These are calibrated
thresholds, not an assumption that formal zero-speed capture vanishes.
The equal-mass capture and transport curves recover the established
calculation to 0.74 per cent and 5.96 per cent in the accepted integrated
log-speed metric.  Nested 4096-versus-8192-particle comparisons pass at
\(q=10^{-2},3\times10^{-3},10^{-3},10^{-4}\); the largest accepted
integrated difference is 18.4 per cent and the largest capture-cutoff shift
is \(|\Delta\ln x|=0.0923\), both below the declared 0.25 limits.
Positive neighbouring values are interpolated logarithmically in speed and
mass ratio; an interval touching a physical zero uses a non-negative linear
transition.  The ordering with \(q\le1\) makes exchange symmetry explicit.
Define the geometric Plummer area per unit mass and the independent
dimensionless collision strength by
\[
 A_{\rm P}\equiv\frac{\pi a^2}{M},
 \qquad S\equiv\code{subhalo_collision_strength}.
\]
The equal-mass channel cross sections per unit mass are the special case
\[
 \frac{\Sigma_j}{M}
 =S A_{\rm P}S_j(x;1),
 \qquad j\in\{\mathrm{cap},\mathrm{scatt}\}.
\]
The runner sets \(A_{\rm P}\) with
\code{subhalo_plummer_area_per_mass}; changing \(S\) multiplies both outcome
channels without changing the homologous subhalo radii.

The analytic limiting behavior is
\[
 S_{\rm cap}(x;q)\simeq\frac{C_{\rm cap}(q)}{x^2}
 \quad\text{when all direct passages capture},
\qquad
 S_{\rm scatt}(x;q)\simeq\frac{C_{\rm T}(q)}{x^4}
 \quad\text{at high }x.
\]
For \(q=1\), \(C_{\rm cap}\simeq52.8\) and
\(C_{\rm T}\simeq168\).  Both coefficients are fixed by the encounter
calculation rather than tuned independently.  The capture cutoff is strongly
mass-ratio dependent: its impulsive small-\(q\) scale decreases as
\(x_{\rm cap,max}\propto q^{1/4}\).  Thus a low-mass object must not inherit
the equal-mass capture curve.

The calibrated domain begins at \(x=0.05\).  The value configured by
\path{subhalo_low_speed_regularization} must be at least this large.  Below
that floor VENUS holds both channel areas at their floor values; this prevents
an unresolved focusing divergence from controlling the Monte--Carlo rate.
At large speed capture is zero and transport scattering uses its asymptotic
\(x^{-4}\) continuation.  Ratios below the smallest calibrated \(q\) use the
verified extreme-mass-ratio row, where fixed nonzero \(x\) capture has
approached zero.

\subsubsection{Unequal masses: physical pair area and whole-object outcomes}

Every marker in this encounter operator is one complete unresolved subhalo.
It is not divided into a colliding packet and an uncollided remainder.  Let
the two masses be \(M_i\) and \(M_j\).  The configured Plummer
area-per-mass \(A_{\rm P}\) implies homologous individual radii
\[
 a_i=\sqrt{\frac{A_{\rm P}M_i}{\pi}},\qquad
 a_j=\sqrt{\frac{A_{\rm P}M_j}{\pi}}.
\]
The physical area scale used by both outcome channels is
\[
 A_{ij}=\frac{A_{\rm P}}{4}
        \left(\sqrt{M_i}+\sqrt{M_j}\right)^2.              \tag{PIC31a}
\]
It is symmetric under \(i\leftrightarrow j\) and becomes \(A_{\rm P}M\)
for two equal masses \(M\).  The pair velocity scale is
\[
 s_{ij}^2=
 \frac{\sqrt{M_i}\,s_i^2+\sqrt{M_j}\,s_j^2}
      {\sqrt{M_i}+\sqrt{M_j}},                            \tag{PIC31b}
\]
where \(s_i\) and \(s_j\) are the stored internal dispersions.  Hence
\[
 \Sigma_{{\rm scatt},ij}=S A_{ij}S_{\rm scatt}(u/s_{ij};q_{ij}),
 \qquad
 \Sigma_{{\rm cap},ij}=S A_{ij}S_{\rm cap}(u/s_{ij};q_{ij}),     \tag{PIC31c}
\]
This recovers the equal-mass boundary at \(q_{ij}=1\), but retains the
separately derived unequal-mass response.
For arbitrary evolved internal dispersions, its unregularized asymptotes are
\[
 \Sigma_{\rm scatt}\simeq
 C_{\rm T}(q_{ij})\,S A_{ij}s_{ij}^{4}u^{-4},
 \qquad
 \Sigma_{\rm cap}\simeq
 C_{\rm cap}(q_{ij})\,S A_{ij}s_{ij}^{2}u^{-2}.          \tag{PIC31ca}
\]
The mass-only statements
\(\Sigma_{\rm scatt}\propto(M_i+M_j)^2u^{-4}\) and
\(\Sigma_{\rm cap}\propto(M_i+M_j)(a_i+a_j)u^{-2}\) follow only while the
objects remain a homologous virial family with
\(s_i^2=C_{\rm vir}\sqrt{M_i}\), where the common constant \(C_{\rm vir}\)
sets the mass--dispersion relation.  Capture heating evolves \(s_i\)
independently, so equation (PIC31ca), not the conditional mass-only scaling,
is the general law used after mergers.

There is an important structural qualification.  The calibrated
\(S_X(x;q)\) rows describe unheated equilibrium Plummer objects with
\(a\propto\sqrt M\).  After capture, VENUS inserts the evolved \(s_i\) in
equation (PIC31b) but retains the homologous radius law.  This is a stated
post-merger closure.  A completely first-principles heated-remnant table would
also require the two dimensionless virial-state parameters
\(\chi_i=s_i^2a_i/(GM_i)\), or an explicit radius-expansion/disruption model.

For collision-cell volume \(V_c\), define
\[
 K_{ij}=u\left(\Sigma_{{\rm scatt},ij}
                    +\Sigma_{{\rm cap},ij}\right).
\]
VENUS does not apply one worst-case cell bound to every pair.  First define
the two generated dimensionless table bounds
\[
 D=D_{\rm scatt}+D_{\rm cap},\qquad
 D_X\ge\max_{\substack{x\ge0\\0<q\le1}}xS_X(x;q).         \tag{PIC31d}
\]
The generator evaluates a conservative endpoint bound for every interpolation
interval, so \(D\) also covers values between stored nodes.
Because the positive weighted root-mean-square in equation (PIC31b) obeys
\(s_{ij}\le\max(s_i,s_j)\), every physical kernel has the rigorous
pair-specific upper bound
\[
 B_{ij}=
 \frac{S A_{\rm P}D}{4}\,
 \max(s_i,s_j)
 \left(\sqrt{M_i}+\sqrt{M_j}\right)^2
 \ge K_{ij},                                               \tag{PIC31e}
\]
where \(S\) is \code{subhalo_collision_strength}.  If
\(B_c=\sum_{i<j}B_{ij}\), the no-time-counter step draws an unordered pair
with probability \(q_{ij}=B_{ij}/B_c\), and proposes
\[
 \mathbb E[N_{\rm trial}]=\frac{B_c\,\delta t}{V_c},
 \qquad
 P_{{\rm accept},ij}=\frac{K_{ij}}{B_{ij}}.                \tag{PIC31f}
\]
Thus the expected accepted rate of that particular pair is exactly
\[
 \frac{B_c\delta t}{V_c}\,
 \frac{B_{ij}}{B_c}\,\frac{K_{ij}}{B_{ij}}
 =\frac{K_{ij}\delta t}{V_c}.
\]
The proposal therefore changes only efficiency, not collision physics.
It avoids making every small, cold pair inherit the largest mass and largest
internal dispersion in the cell.  The density implicit in this
physical-object rate is number density, not mass density.  After acceptance,
the capture probability is
\[
 P_{\rm cap}
 =\frac{\Sigma_{{\rm cap},ij}}
 {\Sigma_{{\rm scatt},ij}+\Sigma_{{\rm cap},ij}}.         \tag{PIC31g}
\]
Each input marker can participate in at most one accepted event per internal
substep.

\paragraph{Elastic scattering.}
Define
\[
 M_{ij}=M_i+M_j,\qquad
 \mu_{ij}=\frac{M_iM_j}{M_{ij}},\qquad
 \bm V=\frac{M_i\bm v_i+M_j\bm v_j}{M_{ij}}.              \tag{PIC32P}
\]
For sampled outgoing relative direction \(\widehat{\bm n}\),
\[
 \bm v_i'=\bm V+\frac{M_j}{M_{ij}}u\widehat{\bm n},
 \qquad
 \bm v_j'=\bm V-\frac{M_i}{M_{ij}}u\widehat{\bm n}.       \tag{PIC33P}
\]
Both complete masses and both internal energies are unchanged.  This map
conserves \(M_i\bm v_i+M_j\bm v_j\) and mass-weighted kinetic energy.  It
creates neither a child marker nor a heavy remainder.

\paragraph{Inelastic capture.}
Capture creates exactly one remnant:
\[
 M'=M_i+M_j,\qquad \bm v'=\bm V,\qquad
 r'=\frac{M_ir_i+M_jr_j}{M_i+M_j}.                        \tag{PIC34P}
\]
The kinetic energy of the relative motion in the pair centre-of-mass frame is
\[
 \Delta K_{\rm rel}=\frac12\mu_{ij}u^2,                   \tag{PIC35P}
\]
and becomes internal energy,
\[
 E_{\rm int}'=E_{{\rm int},i}+E_{{\rm int},j}
              +\frac12\mu_{ij}u^2,\qquad
 s'^2=\frac{2E_{\rm int}'}{3M_{ij}}
 =\frac{M_i s_i^2+M_j s_j^2}{M_{ij}}
  +\frac{\mu_{ij}u^2}{3M_{ij}}.                          \tag{PIC36}
\]
The last positive term is the increase relative to the mass-weighted
progenitor variance.  Therefore capture always heats the remnant relative to
the no-relative-motion mixture, but \(s'\) need not exceed the hotter of
\(s_i\) and \(s_j\).
The second marker is removed during deterministic compaction.  There is no
remainder and therefore no repeated remnant splitting.  Mass, vector
momentum, and centre-of-mass kinetic plus internal energy close to roundoff.
The pair residual does not include gravitational potential energy: replacing
two nearby radii by the mass-weighted radius in equation (PIC34P) can change
the self-consistent potential after redeposition.  That change appears in the
full-system energy history and must converge away with collision-cell
locality and radial resolution.

The source-level structure is:
\codesource{WPIC subhalo collision operator}
{\path{cpu/include/ev/weighted_vlasov_pic_subhalo_collisions.hpp}}
{Applies the pair-specific majorant and then either scatters two complete
 Plummer objects or replaces a captured pair by one energy-carrying remnant}
\begin{lstlisting}[language=C++,basicstyle=\ttfamily\scriptsize]
const double M = first.mass + second.mass;
const double mu = first.mass * second.mass / M;
const Vec3 V = (first.mass * v1 + second.mass * v2) / M;
const double kernel = u * (sigma_scatt + sigma_capture);
const double pair_bound =
    pair_u_sigma_majorant(first, second, options);

if (uniform() < kernel / pair_bound) {
  if (uniform() < sigma_capture / (sigma_scatt + sigma_capture)) {
    first.mass = M;
    first.velocity = V;
    first.internal_energy =
        Eint_first + Eint_second + 0.5 * mu * u * u;
    first.internal_velocity_dispersion =
        sqrt(2 * first.internal_energy / (3 * first.mass));
    remove(second);
  } else {
    first.velocity  = V + (second.mass / M) * u_after;
    second.velocity = V - (first.mass  / M) * u_after;
  }
}
\end{lstlisting}
Cells are independent.  Counter-based random numbers are indexed by epoch,
cell, substep, trial, and random dimension, so thread scheduling cannot
change the physical random stream.

\paragraph{Encounters at the coordinate origin.}
Rotational symmetry makes the velocity distribution at \(r=0\) isotropic.
Every finite velocity there has zero angular momentum.  The reduced marker
therefore stores the full speed as \(w\), with \(\eta=0\).  This convention
also gives the outgoing radial speed when the orbit leaves the origin.
For an encounter, VENUS samples an isotropic Cartesian direction at that
speed.  The pair rate and its accepted outcome use the same sampled vectors.
Scattering and capture conserve momentum in this temporary Cartesian frame.
The outgoing speed is then stored in the marker, preserving its full kinetic
energy:
\codesource{Storing a subhalo velocity at the origin}
{\path{cpu/include/ev/weighted_vlasov_pic_subhalo_collisions.hpp}}
{Stores the magnitude of the Cartesian velocity and zero angular momentum
 at the coordinate origin}
\begin{lstlisting}[language=C++,basicstyle=\ttfamily\scriptsize]
marker.w = std::hypot(velocity[0],
                      std::hypot(velocity[1], velocity[2]));
marker.eta = 0.0;
\end{lstlisting}
The collision-cell dispersion diagnostic uses the isotropic moments relative
to any chosen reference axis at the origin:
\[
 \langle v_r\rangle=0,\qquad
 \langle v_r^2\rangle=\frac{w^2}{3},\qquad
 \langle v_t^2\rangle=\frac{2w^2}{3}. \tag{PIC36a}
\]
Here \(v_r\) is the component along that axis and \(v_t\) is the magnitude
of the two perpendicular components.  The stored \(w\) at the origin is the
speed of radial departure.  Positive-radius encounters retain the usual
radial and tangential velocity components.  A marker at \(r=0\) with
nonzero \(\eta\) is rejected because it has no finite Cartesian velocity.
The CPU and GPU subhalo operators use this same origin convention.

\subsubsection{Collision lanes, pair-majorant proposals, and adaptive time}

A physical radial collision cell can be divided into
\code{subhalo_collision_lanes} deterministic lanes.  Lane membership is an
epoch-hashed stratification of the markers, not an additional radial
coordinate.  Each lane is assigned volume \(V_c/n_{\rm lane}\); therefore
summing its no-time-counter rates gives the original cell rate.  Lanes expose
independent work and reduce the number of candidates in one very populous
cell while preserving the physical collision volume.

The option names \code{maximum_weight_rejection} and
\code{maximum_weight_importance} remain readable for command-line and restart
compatibility.  Both select the same exact pair-majorant proposal in
equations (PIC31e)--(PIC31f).  The runner defaults to
\code{maximum_weight_importance}.  The word ``importance'' means that pairs with a larger rigorous rate bound
are proposed more frequently.  The physical marker masses retain their
stored values throughout this proposal step.

The proposal is constructed without storing all \(N_c(N_c-1)/2\) pairs.
VENUS stably sorts an index array by increasing internal dispersion \(s_i\),
using marker identifier to break ties.  The marker array itself is not
reordered.  For a sorted pair \(i<j\), \(\max(s_i,s_j)=s_j\), and the total
bound carried by upper partner \(j\) is
\[
 G_j=C s_j\sum_{i<j}
 \left(M_i+2\sqrt{M_iM_j}+M_j\right),\qquad
 C=\frac{S A_{\rm P}D}{4}.                                \tag{PIC31h}
\]
Prefix sums of \(M_i\) and \(\sqrt{M_i}\) evaluate every \(G_j\) in one pass.
A cumulative sum of the \(G_j\) first selects \(j\).  Conditional on \(j\),
the lower partner is drawn from three positive pieces proportional to
\(M_i\), \(2\sqrt{M_j}\sqrt{M_i}\), and \(M_j\).  Prefix sums select the first
two pieces; the last is uniform over \(i<j\).  Consequently the exact
probability of selecting \((i,j)\) is \(B_{ij}/B_c\), while setup costs
\(O(N_c\log N_c)\) for sorting and \(O(N_c)\) storage rather than
\(O(N_c^2)\).

The pair proposal and its majorant depend on marker mass and internal
dispersion.  Elastic scattering changes neither quantity, so its exact
proposal is reusable.  Capture changes both and marks only that radial cell
dirty; VENUS removes the consumed object and rebuilds that cell before its
next local substep.  Reusing a pre-capture bound could give an acceptance
probability above one and bias the event rate.  The substep is controlled by
the largest bound on the trial frequency
of any one marker,
\[
 \nu_{B,k}=\frac{1}{V_c}\sum_{j\ne k}B_{kj},\qquad
 \delta t\le
 \frac{\code{subhalo_collision_cfl}}{\max_k\nu_{B,k}}.   \tag{PIC31i}
\]
This is the relevant clock because one marker is allowed at most one accepted
event in a substep.  Prefix and suffix sums evaluate all \(\nu_{B,k}\) in
linear time after sorting.  It is much tighter than multiplying a cell-wide
worst-pair bound by the full population, while remaining conservative.

Reference mode uses one collision clock for the entire radial mesh,
\[
 \delta t_{\rm fixed}=\min\left(
 \Delta t_{\rm op}-t,
 \frac{C_{\rm coll}}{\max_c\nu_{B,c}}\right).           \tag{PIC31i-a}
\]
With \code{--subhalo-adaptive-cell-timestep}, cell \(c\) owns \(t_c\) and
uses
\[
 \delta t_c=\min\left(
 \Delta t_{\rm op}-t_c,
 \frac{C_{\rm coll}}{\nu_{B,c}}\right).                 \tag{PIC31i-b}
\]
Orbital drift is paused during the collision operator.  A captured remnant
is placed at the mass-weighted average of the parent radii, so it remains in
their radial cell.  Every cell therefore advances through the complete
interval independently.  One worker carries out its successive local steps,
rebuilding the proposal after a capture.  OpenMP or the portable thread
backend distributes these cell tasks among workers.  The random key contains
the interval epoch, cell index, local substep, and trial index, so scheduling
preserves every draw.  Diagnostics accumulate in each cell and are then
summed in increasing cell order.  \code{cell_visits} records the sum of local
work; \code{substeps} records the maximum local count.

The collision operator updates the in-memory state as it works.  If a call
fails partway through, continuation must restore an accepted checkpoint.
The driver propagates that failure before resetting event clocks or writing
a checkpoint, preserving the accepted event ordering.

\paragraph{What a low proposal acceptance means.}
For a proposed pair, the physical rate kernel \(K_{ij}\) is accepted with
\[
 P_{{\rm acc},ij}=\frac{K_{ij}}{B_{ij}},\qquad
 0\leq P_{{\rm acc},ij}\leq 1 .                         \tag{PIC31j}
\]
Here \(B_{ij}\) is the rigorous pair majorant used to draw the proposal.  The
accepted fraction of all proposals is therefore an efficiency diagnostic,
not a new physical probability that must exceed an arbitrary percentage.  A
broad unequal-mass or unequal-dispersion population can make \(B_{ij}\)
deliberately conservative for many pairs and drive the aggregate acceptance
below two percent without biasing the accepted event distribution.  Such a
run is expensive, but it is physically valid provided the bound is never
violated, both requested event channels are sampled, the pair invariants and
capture-heating ledger close, the cadence comparison converges, and
\(\Delta r/\lambda_{\rm mfp}<1\).  The runtime diagnostics consequently
report low acceptance as a performance warning rather than rejecting an
otherwise resolved calculation.

\subsubsection{Unresolved origin cell and positive marker compression}

Let \(r_{\min}\) be the first resolved radial mesh face.  The analytic
power-law initializer does not continue its cusp to \(r=0\); it fills
\(0\leq r<r_{\min}\) with a uniform origin core at
\[
 \rho_{\rm core}=\rho(r_{\min}),
\]
where \(\rho(r_{\min})\) is the density at that first resolved face.  Active
markers in this interval contribute to the hidden-core Poisson mass ledger and
remain physical collision partners.

For collisions, let
\[
 I=\code{subhalo_inner_collision_cells}.
\]
VENUS divides the origin-core volume into \(I\) equal-volume physical shells.
Their faces are
\[
 r_a=r_{\min}\left(\frac{a}{I}\right)^{1/3},
 \qquad a=0,\ldots,I,                                    \tag{PIC31k}
\]
so every shell has volume \(4\pi r_{\min}^3/(3I)\).  The central shell is also
the widest:
\[
 \Delta r_{\rm core,max}=r_1-r_0
 =r_{\min}I^{-1/3}.                                     \tag{PIC31l}
\]
Equation (PIC31l) is a statement about the \emph{geometric width only}.
Eight times as many physical inner cells halve that width, but they do not
hold the mean free path fixed.  For the velocity-dependent capture kernel,
each shell must recompute
\(\lambda_{\rm mfp}=\langle u\rangle/
[n_{\rm partner}\langle u\Sigma(u)\rangle]\) from its own marker
distribution.  Resolving a colder central population can strengthen
gravitational focusing and shorten \(\lambda_{\rm mfp}\) faster than
\(\Delta r\) shrinks.  Increasing \(I\) therefore does \emph{not} guarantee a
smaller \(\Delta r/\lambda_{\rm mfp}\).  In contrast, increasing
\code{subhalo_collision_lanes} only supplies independent Monte--Carlo work
inside each physical shell and does not change equation (PIC31l).

The default \(I=1\) reproduces the Poisson solver's single uniform hidden-core
control volume.  Choosing \(I>1\) changes the physical collision
coarse-graining, so it requires a direct per-cell comparison of both
\(\Delta r\) and the newly evaluated \(\lambda_{\rm mfp}\).  The accepted
value must be frozen with \(r_{\min}\), the radial mesh, and the lane count in
the immutable run contract.  It does not place a reflecting wall at
\(r_{\min}\).

\paragraph{Measured origin-core counterexample.}
For the registered \(t=0\), \(\gamma_i=2\) diagnostic with
\(N_r=256\), \(r_{\min}=0.1\), and 8192 initial markers per radial cell, the
maximum locality ratios were
\[
\begin{array}{c|rrrr}
 I & 1 & 2 & 8 & 64\\ \hline
 \max_c(\Delta r_c/\lambda_{{\rm mfp},c}),\ S=1
   &1.20889&1.42762&1.85613&1.96592\\
 \max_c(\Delta r_c/\lambda_{{\rm mfp},c}),\ S=3
   &3.62667&4.28286&5.56839&5.89776
\end{array}                                               \tag{PIC31m}
\]
where \(S\) is the multiplicative encounter-strength parameter.  The
strength-three values are exactly three times the strength-one values for
this kernel and fixed initial state.  Physical cell zero is worst in every
case.  Its lower resolved orbital speed increases the focused capture area,
so its recomputed mean free path decreases more rapidly than the
\(I^{-1/3}\) width.  None of these cases satisfies a unit locality bound.
The lesson is general: select \(I\) from the measured
\(\Delta r/\lambda_{\rm mfp}\), not from geometric width alone, and repeat
that measurement as the distribution evolves.

\paragraph{Why moving the origin inward is not a cure for a singular cusp.}
A second registered \(t=0\) diagnostic held the exterior logarithmic spacing
approximately fixed while moving \(r_{\min}\) inward.  For the same
\(\gamma_i=2\) distribution and \(I=1\), it measured
\[
\begin{array}{c|rrr}
r_{\min} & 0.05 & 0.02 & 0.01\\ \hline
\max_c(\Delta r_c/\lambda_{{\rm mfp},c}),\ S=1
  &2.4082&6.14296&12.4339\\
\max_c(\Delta r_c/\lambda_{{\rm mfp},c}),\ S=3
  &7.22458&18.4289&37.3018
\end{array}                                               \tag{PIC31n}
\]
The strength-three row is again three times the strength-one row.  Over this
range the hidden-cell mean free path falls approximately as
\(\lambda_{\rm mfp}\propto r_{\min}^{\,2}\), whereas its width falls only as
\(\Delta r\propto r_{\min}\).  Consequently
\(\Delta r/\lambda_{\rm mfp}\propto r_{\min}^{-1}\): moving the numerical
origin inward makes the locality error larger.  Increasing only \(N_r\) does
not change the hidden-cell width, and increasing \(I\) already failed in
equation (PIC31m).  Thus a collisional \(r^{-2}\) cusp continued formally to
the origin has no converged hidden-cell limit for this focused encounter law.
It requires a physically specified finite-density core, an inner boundary or
sink, or another short-mean-free-path treatment.  VENUS therefore stops such
a case instead of disguising the unresolved centre with more markers or a
weaker locality threshold.

Whole-object capture reduces the marker count by one; elastic scattering
leaves it unchanged.  Optional
population control calls the positive local quadrature reduction in
\path{canonical_open_carrier_compression.hpp}.  Within local phase subcells it
retains a subset of original phase space points and solves for nonnegative
weights that preserve the selected deposited moments.  For a marker in the
origin handoff, the hidden-carrier contribution \(h_p\) and every resolved
face-cloud contribution \(\widetilde a_{pi}\) in equation (PIC8c) are separate
compression constraints.  Compression therefore cannot undo the smooth
gravity-source handoff by changing the retained markers' distribution within
that interval.  A reduction is rejected
if its relative moment residual exceeds
\code{subhalo_compression_tolerance} or if the fractional reduction is below
\code{subhalo_compression_minimum_reduction_fraction}.  Each retained packet
also preserves its specific internal energy \(E_{\rm int}/m\).

Compression begins only when the total count reaches
\code{subhalo_compression_trigger_markers}; the run may retain at least
\code{subhalo_compression_minimum_markers}.  The source-cell preservation flag
forbids a replacement packet from crossing its original phase space source
cell.  Compression is a controlled approximation to the sampled distribution,
so its counts, residuals, timings, and compressed/uncompressed profile
differences must be included in the physical convergence study.

\subsubsection{Using unequal weights to overcome the SIDM marker barrier}
\label{sec:wpic-unequal-sidm-path}

Section~\ref{sec:sphdsmc-marker-scaling} shows why a globally equal marker
mass becomes prohibitive during SIDM core collapse:
\[
 N_{\rm req}\propto
 \frac{(k-1)^3}{k^2}
 \left(\frac{\sigma}{m}\right)^3\rho_{\rm target}^2
 \ \longrightarrow\
 k\left(\frac{\sigma}{m}\right)^3\rho_{\rm target}^2
 \quad(k\gg1).
\]
The WPIC state has the representation needed to avoid this global cost.
Collisionless characteristics carry an individual numerical weight \(m_i\),
and mass deposition sums those weights before solving the Poisson equation.
The outer halo can therefore be represented by relatively heavy markers while
the collapsing core is represented by much lighter markers.

For a local core marker mass \(m_{\rm core}\), the central locality estimate is
\[
 \frac{\Delta r}{\lambda_{\rm mfp}}
 \simeq
 \frac{k-1}{k}\,
 \frac{\sigma}{m}
 \left(\frac{3k m_{\rm core}}{4\pi}\right)^{1/3}
 \rho_{\rm core}^{2/3}.                                  \tag{PIC34cW}
\]
Here \(k\) is the local collision-cell population, so \((k-1)/k\) excludes
the selected marker itself from its partner density.  It tends to unity in
the continuum limit.
To enforce a requested maximum ratio \(\chi_{\rm target}\), the core weight
must satisfy
\[
 m_{\rm core}\lesssim
 \frac{4\pi}{3k}\,
 \left(\frac{k}{k-1}\right)^3
 \frac{\chi_{\rm target}^3}
 {(\sigma/m)^3\rho_{\rm target}^2}.                       \tag{PIC34dW}
\]
This is the unequal-weight form of the equal-mass requirement: it constrains
the marker mass only where the density is high, rather than forcing every
marker in the outer halo to have the same very small mass.

Let a baseline calculation have \(N_0\) equal markers of mass \(m_0=M/N_0\).
If the refined core needs \(m_{\rm core}=m_0/F\) and contains a fraction
\(f_c=M_c/M\) of the represented mass, retaining the original outer
resolution gives approximately
\[
 N_{\rm adaptive}
 \simeq N_0\left[1+f_c(F-1)\right],                       \tag{PIC34e}
\]
whereas uniform refinement would require \(N_{\rm uniform}=FN_0\).
Equation (PIC34e) quantifies the computational advantage: when only a small
mass fraction requires the finest weight, adaptive marker counts can be far
below the uniformly refined count \(FN_0\).  The refined region must include
the transition zone through which collisions conduct heat, not only the
instantaneous density core.  Use equations (S14c) and (PIC34dW) with the
selected shell population \(k\), then recompute the locality ratio from the
evolving velocity distribution throughout the calculation.

\paragraph{Implemented unequal-mass operations.}
The general \code{WeightedPICMarker} stores an independent represented mass,
and the collisionless transport, field deposition, restart, and diagnostics
preserve it.  The ordinary material-SIDM and gravitational WPIC kernels, the
Frank/SphDSMC restart importer, and the invariant-preserving phase remap
accept arbitrary positive unequal masses.  For a configured material law
\(s(u)=\sigma(u)/m\), the physical symmetric pair area is
\[
 \boxed{\Sigma_{ij}(u)=\frac{m_i+m_j}{2}\,s(u)}.          \tag{PIC34f}
\]
For \(m_i=m_j=m_p\), number density
\(n=\rho/m_p\) gives
\(n\Sigma_{ij}=\rho s\) in the continuum limit, exactly the original
resolution-independent collision frequency.  A finite \(N_s\)-object cell
uses \(n_{\rm partner}=n(N_s-1)/N_s\), as in equation (PIC14c).  For unequal
partners, the collision map uses
\[
 \bm V=\frac{m_i\bm v_i+m_j\bm v_j}{m_i+m_j},\qquad
 \bm v_i'=\bm V+\frac{m_j}{m_i+m_j}\bm u',\qquad
 \bm v_j'=\bm V-\frac{m_i}{m_i+m_j}\bm u'.                \tag{PIC34g}
\]
Both full masses remain unchanged.  The Plummer subhalo operator uses the
whole-object scattering and capture rules in equations
(PIC31a)--(PIC36).  Every collision path accepts positive unequal masses, and
no accepted event creates a split remainder.

\paragraph{Interpretation and adaptive-resolution requirements.}
Equation (PIC34f) treats a marker mass as the mass of one complete
coarse-grained collision partner.  This is the intended closure for the
present unequal-mass code.  Arbitrarily changing marker weights is therefore
not a neutral Monte-Carlo refinement: it changes the represented mass
spectrum unless the initialization or resampling prescription is constructed
to preserve the desired physical mixture.  A production adaptive SIDM
calculation still requires:
\begin{enumerate}
\item monitor \(\Delta r/\lambda_{\rm mfp}\) and generate a lighter,
      well-sampled core population before a collision region becomes nonlocal;
\item place newly sampled core representatives at distinct local phase space
      points; exact duplicates conserve moments but do not improve spatial
      sampling;
\item build collision regions from their number density, pair cross sections,
      represented mass, and geometry, and
      require a sufficient weighted effective population
      \[
       N_{\rm eff}=
       \frac{(\sum_i m_i)^2}{\sum_i m_i^2};
      \]
\item use conservative paired resampling to construct the intended mass
      spectrum while preserving mass, momentum, and kinetic-energy moments;
\item refine the radial gravity mesh with the marker distribution, because
      light core markers cannot recover a core smaller than the Poisson-grid
      resolution;
\item compare adaptive-resampling sequences through the desired core-collapse
      interval, including Maxwellian null behavior, event/remap conservation,
      collision rates, locality, and gravothermal heat transport.
\end{enumerate}
Unequal weights remove the need to refine the entire halo uniformly.  They do
not change the physical mean free path, prevent the core from entering the
short-mean-free-path regime, or remove the shrinking collision and orbital
timesteps at late collapse.

\subsubsection{Subhalo collision diagnostics and locality}

For this operator the fixed radial WPIC mesh cells are collision volumes.
Before computing locality, VENUS reunites the parallel Monte-Carlo lanes that
belong to one physical radial cell.  Lanes accelerate independent pair trials;
they are not separate pieces of matter and therefore cannot alter a physical
mean free path.  With
\(\Sigma_{ij}=\Sigma_{{\rm scatt},ij}+\Sigma_{{\rm cap},ij}\), the code uses
\[
 \lambda_{{\rm mfp},c}=
 \frac{\langle u_{ij}\rangle_c}
 {n_{{\rm partner},c}
  \langle u_{ij}\Sigma_{ij}(u_{ij})\rangle_c},\qquad
 \sigma_{v,c}^{\,2}=\frac{1}{3M_c}\sum_{i\in c}m_i
 \left[(w_i-\overline w_c)^2+\frac{2\eta_i}{r_i^2}\right].
\]
Here \(c\) labels a physical radial cell, \(N_c\) is its marker count,
\(V_c\) is its volume,
\(n_c=N_c/V_c\) is its object-density estimate,
\(n_{{\rm partner},c}=(N_c-1)/V_c\) is the density of other objects available
to one selected marker, \(M_c=\sum_i m_i\), and every other symbol has the
same meaning as in equations (PIC14a)--(PIC14c).  Inner cells use 4096
deterministic pair samples and other cells use 256, or all pairs if fewer
exist.  The code reports
\[
 \frac{\Delta r_c}{\lambda_{{\rm mfp},c}}
\]
and compares it with \code{--subhalo-locality-maximum-ratio}.  A zero
threshold leaves the quantity diagnostic-only.  A positive threshold causes
an atomic checkpoint and status 98 at the first violation; production uses
one, so that no accepted output silently enters
\(\Delta r_c>\lambda_{{\rm mfp},c}\).  The violation is a resolution failure:
the collision cell is wider than the estimated local mean free path.
Decreasing a physical cell width, reducing the unresolved \(r_{\min}\)
region, or changing the physical coarse-graining may help, but every candidate
must be re-evaluated with the evolving velocity distribution.  More markers
improve the statistical estimate but do not by themselves lengthen the
physical mean free path; increasing \(I\) is not guaranteed to reduce the
ratio, as equation (PIC31m) demonstrates.  A numerical cell face must never
be treated as a physical boundary merely to bypass the stop.

\paragraph{Does repeated capture itself cause the locality loss?}
Capture can accelerate central contraction because relative orbital kinetic
energy is transferred into the remnant's internal degrees of freedom rather
than returned to the centre-of-mass orbit.  A denser centre then tends to
increase \(n_{{\rm partner},c}\langle u\Sigma\rangle_c\) and shorten
\(\lambda_{{\rm mfp},c}\).  It is not correct, however, to infer the mean free
path from central mass density alone.  Each capture reduces the marker number
\(N_c\) by one, increases a remnant mass and its pair area, changes the local
relative-speed distribution, and adds the positive heating term in equation
(PIC36) relative to the mass-weighted progenitor variance.  The remnant
dispersion need not exceed that of the hotter progenitor.  These effects enter
the numerator and denominator of the measured expression above with different
signs.  VENUS therefore recomputes the pair average and the mass-weighted
velocity dispersion from the evolving distribution after every collision
application.  A rising central density is a physical explanation for why
locality may become difficult; the actual failure is established only by the
measured \(\Delta r_c/\lambda_{{\rm mfp},c}\), not by capture count or density
alone.

\path{collision_history.csv} records the summed number of substeps and trials
since the preceding history write,
accepted scattering and capture events, represented mass in each channel,
removed markers (the spawned-marker column is zero), orbital and internal-energy changes, total
internal energy, marker count, maximum marker-weight ratio, locality counts,
the number of mean-free-path pair samples, the minimum collision,
dynamical, and potential-change times, the minimum/maximum applied collision
interval, the maximum interval divided by the physical clock measured just
before application, the maximum interval divided by the clock reconstructed
after the application, and the
minimum/maximum evolving orbital velocity dispersion,
minimum/maximum marker and pair internal dispersions, maximum pair momentum
and total-energy residuals, and the compression counts, before/after marker
populations, retained groups, maximum moment residual, and wall time.  These
columns separate a physical transfer of orbital energy into internal energy
from numerical nonconservation and expose every population-control event.
The CSV convention uses \code{min_potential_change_time=0} to record
\(t_\Phi=\infty\): the measured field change stayed below the finite-marker
noise floor.  This stored zero therefore denotes an unresolved change rate
and an infinite potential-change time.
The two interval ratios are respectively
\code{max_interval_over_selection_timescale} and
\code{max_interval_over_retrospective_timescale}.  They are diagnostics,
not conservation laws, but they have different acceptance meanings.  For a
positive \code{--subhalo-collision-timescale-fraction}, the selection ratio is
a hard pre-collision bound: the runner rolls back a provisional orbit step
rather than apply a batch whose complete accumulated age exceeds the
currently measured deadline.  The collision itself can nevertheless shorten
the post-event clock by changing marker masses, pair areas, the velocity
distribution, and the potential.  Consequently the retrospective ratio can
transiently exceed the configured fraction.  That post-event excess requests
a cadence-convergence check; it does not retroactively invalidate an
otherwise time-ordered collision batch.

The substep count is a history-window sum, not the maximum used by one
application of the collision operator.  The per-application
\code{subhalo_collision_maximum_substeps} guard is enforced inside the
operator; reaching the next history record proves that none of the completed
applications hit that guard.

\subsubsection{Four separate timing controls}

Four controls govern the calculation: the orbit step resolves trajectories,
the interval between collision applications resolves transport--collision
coupling, the internal collision substeps resolve event sampling, and a
maximum-substep guard bounds the computational work.  The following
definitions explain the physical timescales and how each control uses them.

\paragraph{Physical collision time.}
In physical radial cell \(c\), let
\[
 {\cal K}_c=\left\langle
 u_{ij}\Sigma_{ij}(u_{ij})\right\rangle_c,
 \qquad
 \nu_{{\rm coll},c}=n_{{\rm partner},c}{\cal K}_c .
\]
Here \(u_{ij}=|\boldsymbol v_i-\boldsymbol v_j|\) is pair relative speed,
\(\Sigma_{ij}=\Sigma_{{\rm scatt},ij}+\Sigma_{{\rm cap},ij}\) is the total
unequal-mass pair area,
\(n_{{\rm partner},c}=(N_c-1)/V_c\) is finite-cell partner density, and the
angle brackets denote the evolving pair sample in that cell.  The mean
waiting time for one selected object to encounter another and the mean free
path are
\[
 t_{{\rm coll},c}=\frac{1}{\nu_{{\rm coll},c}}
 =\frac{1}{n_{{\rm partner},c}\langle u\Sigma\rangle_c},
 \qquad
 \lambda_{{\rm mfp},c}
 =t_{{\rm coll},c}\langle u\rangle_c .                 \tag{PIC37}
\]
Equation (PIC37) shows why the velocity dispersion cannot be omitted for a
velocity-dependent cross section: the evolving distribution changes both
\(\langle u\rangle\) and \(\langle u\Sigma(u)\rangle\).

\paragraph{Dynamical time.}
The physical enclosed marker mass is
\[
 M(<r_i)=M_{\rm hidden}+M_\bullet
          +\sum_{j\leq i}M_{{\rm cell},j}.
\]
Here \(M_{\rm hidden}\) is active marker mass currently inside the first
resolved radial face, \(M_\bullet\) is the configured central point mass, and
\(M_{{\rm cell},j}\) is the deposited physical mass in radial cell \(j\).
The expanding runner normally has \(M_\bullet=0\), but the shared field
routine retains the term.
The corresponding spherical gravitational frequency and time are
\[
 \Omega_{{\rm dyn},i}^2=\frac{G M(<r_i)}{r_i^3},
 \qquad
 t_{{\rm dyn},i}=\Omega_{{\rm dyn},i}^{-1}.             \tag{PIC38}
\]
VENUS uses the minimum positive \(t_{{\rm dyn},i}\) over represented cells.
It uses total enclosed matter in equation (PIC38), not the density contrast
used as the cosmological Poisson source.

\paragraph{Potential-change time.}
Let \(g=-\partial\Phi/\partial r\).  Successive self-consistent Poisson fields
first give the dimensionless finite changes
\[
 \delta_\Phi=
 \frac{\|\Phi^{n+1}-\Phi^n\|_\infty}
      {\|\Phi\|_{\infty,*}},
 \qquad
 \delta_g=
 \frac{\|g^{n+1}-g^n\|_\infty}
      {\|g\|_{\infty,*}},                              \tag{PIC39a}
\]
where
\[
 \|\Phi\|_{\infty,*}=\max_i\max(|\Phi_i^n|,|\Phi_i^{n+1}|),
 \qquad
 \|g\|_{\infty,*}=\max_i\max(|g_i^n|,|g_i^{n+1}|),
\]
so the normalization does not vanish merely because one grid point crosses
zero.

A finite marker representation also produces a small deposition residual
between nearly identical fields.  Let \(a_{pc}\) be the fraction of marker
\(p\)'s mass deposited in mesh cell \(c\), \(m_p\) its represented mass, and
\(M_c=\sum_p a_{pc}m_p\).  The global deposited effective marker number used
by the clock is
\[
 N_{\rm eff}^{\rm glob}
 =\frac{\left(\sum_c M_c\right)^2}
        {\sum_c\sum_p(a_{pc}m_p)^2}.                   \tag{PIC39b}
\]
It decreases when a few unequal-mass markers dominate the field even if the
raw marker count is large.  VENUS defines the resolved-change threshold
\[
 \delta_N=\max\left(4096\,\epsilon_{\rm mach},
                    \frac{0.25}{\sqrt{N_{\rm eff}^{\rm glob}}}\right),
\]
where \(\epsilon_{\rm mach}\) is double-precision machine roundoff.  For
\(X\in\{\Phi,g\}\), the measured change time is
\[
 t_X=
 \begin{cases}
 \displaystyle
 \frac{\Delta t_{\rm evo}}
      {\sqrt{\delta_X^2-\delta_N^2}},&
             \delta_X>\delta_N,\\[6pt]
 \infty,& \delta_X\leq\delta_N ,
 \end{cases}
 \qquad t_\Phi=\min(t_\Phi^{(\Phi)},t_\Phi^{(g)}).      \tag{PIC39c}
\]
The factor \(0.25\) is deliberately permissive: it is one quarter of the
nominal \(N_{\rm eff}^{-1/2}\) sampling scale, so a marginally resolved change
is more likely to shorten the cadence than to be hidden.  Sub-threshold
changes are not converted into a derivative because dividing a deposition
residual by a tiny output-cleanup step would manufacture
\(t_\Phi\rightarrow0\).  This threshold distinguishes resolved field evolution from deposition noise.
Its statistical calibration requires independent realizations.  Marker-number and quiet-start convergence should assess
it; a future refinement may propagate cellwise deposited-mass variance
through the Poisson solve.  Thus \(t_\Phi\) is measured from resolved evolution
of the self-consistent field rather than assumed equal to \(t_{\rm dyn}\).

\paragraph{How often the collision generator is applied.}
Define
\[
 t_{\rm fast}=\begin{cases}
 \min_c\left(t_{{\rm dyn},c},t_\Phi\right),
 &\text{CPU adaptive collision cells},\\
 \min_c\left(t_{{\rm coll},c},t_{{\rm dyn},c},t_\Phi\right),
 &\text{mesh-wide reference and CUDA paths}.
 \end{cases}
\]
The outer collision-operator interval is
\[
 \Delta t_{\rm op}=
 \min\left(\Delta t_{\rm op,max},
            \epsilon_{\rm op}t_{\rm fast}\right).       \tag{PIC40}
\]
\(\Delta t_{\rm op,max}\) is
\code{--subhalo-collision-interval}, and
\(\epsilon_{\rm op}\) is
\code{--subhalo-collision-timescale-fraction}.  The baseline production
template starts with \(\epsilon_{\rm op}=0.1\); a short run that fails the
registered interval-halving comparison must use a smaller value, commonly
0.05, and remains provisional until the next halving agrees within the
profile tolerances.  In the mesh-wide reference and CUDA paths,
\(t_{\rm coll}\) sets the outer deadline when it is shorter than both
\(t_{\rm dyn}\) and \(t_\Phi\).  The CPU
\code{--subhalo-adaptive-cell-timestep} path instead recomputes collision rates
inside each cell and enforces its local collision CFL there; the outer
operator interval is limited by the dynamical and potential-change clocks.
The CPU selection and endpoint checks pass an infinite collision time in this
mode (\path{cpu/src/drivers/expanding_wpic.cpp}, lines 2285--2313 and
2686--2701).  CUDA retains the sampled collision time in its outer minimum
(\path{gpu/src/wpic_production_driver.cu}, lines 1070--1083).
A diagnostic pair scan estimates \(t_{\rm coll}\)
before the first real event without changing markers or consuming the
production random stream.  Version-6 restarts preserve the latest clock
estimates.

Where \(t_{\rm coll}\) enters equation (PIC40), its latest measured value is
held between collision applications.  The evolving pair diagnostic refreshes
it at each application.  The dynamical and noise-resolved potential-change
clocks are re-evaluated at a provisional orbit endpoint, using the same
branch of \(t_{\rm fast}\).  Outer-cadence refinement remains necessary for
each path.

An adaptive KDK step is transactional.  Before the step, VENUS saves every
marker radius \(r_p\) and radial velocity \(w_p\).  It then performs the first
kick and drift, deposits the provisional endpoint distribution, solves
Poisson, and recomputes \(t_{\rm dyn}\) and \(t_\Phi\).  If
\[
 t_{\rm age}+\Delta t_{\rm evo}>
 \min\!\left(\Delta t_{\rm op,max},
             \epsilon_{\rm op}t_{\rm fast,end}\right), \tag{PIC40a}
\]
the endpoint is rejected: \(r_p,w_p\) are restored and a shorter orbit step
is tried.  Here \(t_{\rm age}\) is physical time accumulated since the last
collision-operator application.  A rejected probe advances neither physical
time, accepted step number, collision age, nor the counter-based collision
random stream.  Only after a probe satisfies equation (PIC40a) is the endpoint
kick completed and the state committed.  A second hard check is made
immediately before the collision generator is called.  Therefore
\code{max_interval_over_selection_timescale} cannot exceed
\(\epsilon_{\rm op}\), apart from the stated roundoff tolerance, in an
adaptive accepted run.

The endpoint comparison includes only a few units in the last represented
place.  Let
\[
 D_{\rm end}=
 \min\!\left(\Delta t_{\rm op,max},
             \epsilon_{\rm op}t_{\rm fast,end}\right),
 \qquad
 \delta_D=64\epsilon_{\rm mach}
 \max(1,|t_{\rm age}|,|D_{\rm end}|),                    \tag{PIC40b}
\]
where \(D_{\rm end}\) is the newly measured collision deadline and
\(\epsilon_{\rm mach}\) is the machine relative spacing for the stored
double-precision numbers.  A pre-step collision age is still timely only
when
\[
 t_{\rm age}\leq D_{\rm end}+\delta_D.                   \tag{PIC40c}
\]
The smallest useful retry at coordinate time \(t\) is estimated as
\[
 \delta t_{\rm coord}=128\epsilon_{\rm mach}
                       \max(1,|t|),                      \tag{PIC40d}
\]
where \(\delta t_{\rm coord}\) is the local resolution of the floating-point
time coordinate.  If the provisional endpoint violates equation (PIC40a),
VENUS first restores the saved pre-step marker coordinates.  It normally
retries a shorter interval that remains inside \(D_{\rm end}\).  In the
exceptional case
\[
 D_{\rm end}-t_{\rm age}\leq\delta t_{\rm coord}
 \quad\hbox{and}\quad
 t_{\rm age}\leq D_{\rm end}+\delta_D,                   \tag{PIC40e}
\]
no positive representable KDK step can reach the still-timely event
reliably.  VENUS then applies the collision batch at the rolled-back pre-step
state.  This is a coordinate-resolution-limited, infinitesimally early
application, not permission to apply an overdue operator.

The field or pair distribution produced by the collision itself can shorten
the newly measured clock.  Hence
\code{max_interval_over_retrospective_timescale} may exceed
\(\epsilon_{\rm op}\); it records post-event physics rather than the deadline
known before the batch.  If instead
\(t_{\rm age}>D_{\rm end}+\delta_D\), the pre-step state is genuinely overdue:
safe-early application is forbidden.  When rollback refinement reaches
\(\delta t_{\rm coord}\) in that state, VENUS fails closed, checkpoints the
pre-collision state without sampling the batch, and exits with status 97.
The same fail-closed result is used if the finite rollback-probe guard is
exhausted before any admissible ordering is established.  Such a checkpoint
can already contain a
collision-event age longer than the newly discovered deadline.  Merely
restarting that same state with a smaller cadence fraction cannot restore the
missing earlier operator application.  A valid recovery must begin from an
earlier checkpoint whose saved age satisfies the new deadline (normally a
collision boundary), or from a fresh run with the refined cadence.

\paragraph{Ordinary evolution timestep.}
Between collision deadlines, KDK Vlasov evolution has its own timestep.  In
mesh cell \(i\), the implemented frequency estimate is
\[
 \Omega_{{\rm evo},i}^2=
 \frac{|\Phi_r(r_i)|}{r_i}
 +\frac{\langle v_t^2\rangle_i}{r_i^2}+H^2 ,
\]
where \(\langle v_t^2\rangle_i\) is the deposited tangential second moment
and \(H\) is the expansion rate.  The minimum is taken only over represented
cells satisfying
\[
 M_{{\rm cell},i}>10^{-12}M_{\rm mesh},
\]
where \(M_{\rm mesh}\) is the total mass deposited on the resolved mesh.
This prevents an almost empty outer cell from setting the global orbit step.
For a collisional run, therefore,
\[
 \Delta t_{\rm evo}=
 \min\left[\Delta t_{\rm evo,max},
 \min_i\frac{\epsilon_{\rm dyn}}{\Omega_{{\rm evo},i}},
 \Delta t_{\rm next\ output},
 \Delta t_{\rm next\ collision}\right].                \tag{PIC41}
\]
\code{--maximum-timestep} supplies \(\Delta t_{\rm evo,max}\), while
\code{--dynamical-fraction} supplies \(\epsilon_{\rm dyn}\).  The last Vlasov
step is shortened to land exactly on an output or collision deadline.  In an
adaptive collisional run this is still only the proposed KDK interval; the
endpoint audit in equation (PIC40a) may roll it back and shorten it again.
For a collisionless WPIC run VENUS deliberately does not apply the cellwise
dynamical limiter: it uses the fixed
\(\Delta t_{\rm evo,max}\), shortened only to land on an output or the final
time.  Fixed-step time-symmetric KDK gives bounded energy error in a static
potential; a varying step would destroy that symmetry even when the force is
constant in time.

\paragraph{Internal collision CFL.}
One application covering \(\Delta t_{\rm op}\) is further divided.  In
reference mode, if \(\nu_{B,k}\) is the safe pair-majorant trial-frequency
bound for marker \(k\), VENUS takes
\[
 \delta t_{\rm coll}=
 \min\left(\Delta t_{\rm op}-t_{\rm elapsed},
 \frac{\epsilon_{\rm coll}}{\max_k\nu_{B,k}}\right),    \tag{PIC42}
\]
where \(\epsilon_{\rm coll}\) is
\code{--subhalo-collision-cfl}.  The majorant is rebuilt after capture changes
the local population.  Adaptive-cell mode replaces the mesh maximum by the
local equation (PIC31i-b).  The maximum-substep option is a safety guard, not a
replacement for equations (PIC40)--(PIC42).

\paragraph{Administrative maximum-substep guard.}
\code{--subhalo-collision-maximum-substeps} limits only how many
equation-(PIC42) substeps one operator application may execute before the
runner deliberately aborts.  It neither chooses \(\delta t_{\rm coll}\) nor
changes an NTC probability.  The executable default of 256 catches accidental
runaway work in small runs.  Resolved high-rate unequal-mass evolution can
legitimately require many more CFL-selected substeps; a practical initial
guard is 16384.  An especially steep or strongly collisional interval can
still exhaust 16384; retrying that same case with a larger guard, for example
32768 or 65536, is legitimate
when the collision CFL, locality, and invariants remain unchanged.  Raising
the guard does not change physical accuracy or the collision sequence
provided the lower value would otherwise only have stopped the same
application.  It merely permits the CFL-controlled interval to finish.
Conversely, hitting the guard produces an error rather than accepting a
partially integrated collision interval.

These intervals need not have a permanent ordering: a very dense cell can
have \(\delta t_{\rm coll}<\Delta t_{\rm evo}\), while a weakly collisional
halo can take many Vlasov steps before one outer application.  The outer
frequency also requires empirical refinement.  If \({\cal V}\) denotes Vlasov
flow and \({\cal C}\) the collision generator, sequential splitting obeys
\[
 e^{\Delta t{\cal C}}e^{\Delta t{\cal V}}
 =e^{\Delta t({\cal V}+{\cal C})
   +\frac12\Delta t^2[{\cal C},{\cal V}]
   +O(\Delta t^3)} .                                   \tag{PIC43}
\]
The commutator \([{\cal C},{\cal V}]\) expresses that collision rates change
as gravity moves the distribution.  The cadence test therefore repeats short runs with the same random seed and \(\Delta t_{\rm op,max}\) halved and compares density,
logarithmic slope, capture and scattering counts, conservation, and
\(\Delta r/\lambda_{\rm mfp}\).  Independent-seed ensembles must also be
compared: the mean density and slope profiles from the two cadences should
agree within their measured Monte--Carlo scatter.  Agreement indicates that
the splitting term in equation (PIC43) is small; merely staying below the
internal collision CFL is not sufficient.  Irrespective of either timing
ratio, collision-cell locality and mass, momentum, and total-energy
conservation remain hard acceptance gates.

\paragraph{Cadence convergence requirements.}
No outer collision cadence is a universal property of the physical model.
For every marker loading, cusp, cross section, and radial mesh, perform a
refinement with the same random seed of the pair
\((\Delta t_{\rm op,max},\epsilon_{\rm op})\) before examining its result.
Halving both is a conservative joint refinement; halving only one is
informative only after the runtime diagnostics show that control is active.
Require the intended bounds to be reached in practice, then compare density,
logarithmic slope,
scattering and capture counts, conservation ledgers, and
\(\Delta r/\lambda_{\rm mfp}\).  Repeat the comparison over independent paired
collision seeds so that a cadence effect can be distinguished from ordinary
Monte--Carlo scatter.  Do not interpret a profile until the chosen differences
are smaller than the scientific error budget.  The time ordering is invalid
if any accepted pre-event interval violates
\[
  \code{max_interval_over_selection_timescale}
  \leq \epsilon_{\rm op}(1+2\times10^{-12}).
\]

\subsubsection{Running expanding WPIC with scattering and capture}

Build the dedicated expanding task frontend with:
\codesource{\code{venus_cpu_subhalo_wpic}}
{\path{cpu/src/apps/subhalo/expanding_wpic.cpp};
 \path{cpu/Makefile}}
{Builds the fail-closed expanding subhalo WPIC frontend}
\begin{lstlisting}
cd VENUS/cpu
make venus_cpu_subhalo_wpic
\end{lstlisting}
A collisionless control uses the Vlasov--Poisson frontend with zero strength:
\codesource{\code{venus_cpu_vlasov_poisson_wpic_expanding}}
{\path{cpu/src/apps/vlasov_poisson/expanding_wpic.cpp};
 \path{cpu/src/drivers/expanding_wpic.cpp}}
{Runs the same expanding marker and gravity problem with the subhalo collision
 strength set exactly to zero}
\begin{lstlisting}
./venus_cpu_vlasov_poisson_wpic_expanding \
  --gamma 1.5 --end-time 2 \
  --particles-per-radial-cell 2048 \
  --subhalo-collision-strength 0 \
  --threads 32 \
  --output-directory output/wpic_gamma15_collisionless
\end{lstlisting}
The corresponding scattering/capture run uses the subhalo frontend with the same
initial data and gravity settings, and enables the collision block:
\codesource{\code{venus_cpu_subhalo_wpic}}
{\path{cpu/src/apps/subhalo/expanding_wpic.cpp};
 \path{cpu/src/drivers/expanding_wpic.cpp}}
{Runs the paired collisional case with whole-object Plummer scattering and
 capture enabled}
\begin{lstlisting}
./venus_cpu_subhalo_wpic \
  --gamma 1.5 --end-time 2 \
  --particles-per-radial-cell 2048 \
  --subhalo-collision-strength 1 \
  --subhalo-velocity-scale 0.1 \
  --subhalo-plummer-area-per-mass 1 \
  --subhalo-low-speed-regularization 0.05 \
  --subhalo-collision-interval 1 \
  --subhalo-collision-timescale-fraction 0.1 \
  --subhalo-collision-cfl 0.2 \
  --subhalo-collision-maximum-substeps 16384 \
  --threads 32 \
  --output-directory output/wpic_gamma15_subhalo
\end{lstlisting}
\code{--subhalo-collision-interval} is a ceiling on the physical time between
operator applications.  With
\code{--subhalo-collision-timescale-fraction 0.1}, the actual interval is
normally shorter and is selected from the evolving physical clocks derived
below.  It must still be refined until profiles, logarithmic slopes, and
channel histories are insensitive to halving either the ceiling or the
timescale fraction.  The NTC substeps control event probability inside one
application; they do not remove the need to converge the outer interval.
The value 16384 is an administrative production guard, not a timestep or an
accuracy target.

A restart continuation uses:
\codesource{\code{venus_cpu_subhalo_wpic}}
{\path{cpu/src/drivers/expanding_wpic.cpp}}
{Restores the complete marker, collision-sampler, physical-clock, and radial-
 mesh contract before continuing the expanding run}
\begin{lstlisting}
./venus_cpu_subhalo_wpic \
  --gamma 1.5 --end-time 4 \
  --restart output/wpic_gamma15_subhalo/restart_state.bin \
  --subhalo-collision-strength 1 \
  --subhalo-velocity-scale 0.1 \
  --subhalo-plummer-area-per-mass 1 \
  --subhalo-low-speed-regularization 0.05 \
  --subhalo-collision-interval 1 \
  --subhalo-collision-timescale-fraction 0.1 \
  --subhalo-collision-cfl 0.2 \
  --subhalo-collision-maximum-substeps 16384 \
  --threads 32 \
  --output-directory output/wpic_gamma15_subhalo
\end{lstlisting}
The current runner publishes \path{restart_state.bin} atomically.
Collision settings must be kept physically consistent with the original run.

\subsubsection{Weighted-PIC restart format and present limitations}

Weighted-PIC checkpoints use magic \code{VENUSWVPIC1}. Format version 6
stores marker mass and phase coordinates, state,
boundary data, identifier, internal energy, internal velocity dispersion,
initialization sampler, collision sampler, potential and collision ages,
the most recent collision and potential-change timescale estimates, physical
time, and step.  The writer uses a temporary file plus atomic rename and
checks identifier uniqueness. The static Boltzmann event driver writes
version 7, which adds the frozen Poisson field and its evolution-control
record after the marker records. CPU readers accept versions 1--7; older
versions initialize fields that did not exist to their documented defaults.
Trailing bytes, short reads, invalid states, and duplicate identifiers are
errors.

Implementation boundaries of the WPIC modules are:
\begin{itemize}
\item adaptive core refinement and a physics-preserving resampling rule for a
chosen unequal-mass spectrum remain to be implemented; although the phase
remap accepts every positive marker mass, a nearly continuous mass spectrum
usually takes the duplicated-singleton path and still requires the
carrier/cost/refinement preflight described above before science use;
\item the subhalo module uses complete unequal-mass objects and fixed radial
mesh cells rather than adaptive mean-free-path shells;
\item subhalo scattering uses the tabulated momentum-transfer area with an
isotropic unresolved outgoing direction;
\item capture heating is conserved, checkpointed, and changes the later
internal velocity scale, but no Plummer-radius expansion/disruption law is
implemented;
\item collision lanes and the pair-majorant proposal require
full-resolution equivalence comparisons before final science use; the two
pair-proposal option names both select the same strict pair-bound
importance sampler;
\item optional positive compression preserves monitored moments but still
requires compressed/uncompressed population-control convergence;
\item the static event-coupled WPIC wrapper and the expanding subhalo runner
are separate integration paths;
\item locality, marker-weight ratio, collision-interval, radial resolution,
and marker count all require convergence checks.
\end{itemize}

\subsubsection{Equilibrium properties and convergence criteria}

For equal represented masses, a drifting Maxwellian with mean velocity
\(\bm U\) and one-component velocity variance \(\sigma^2\) obeys
\[
 f_M(\bm v_1)f_M(\bm v_2)\propto
 \exp\left[-\frac{|\bm v_1-\bm U|^2+|\bm v_2-\bm U|^2}
 {2\sigma^2}\right]. \tag{PIC35M}
\]
Conservation of pair momentum and kinetic energy in equation (PIC27)
preserves this exponent, so forward and reverse events balance.  For unequal
object masses \(m_1,m_2\), the corresponding equilibrium product contains
\[
 \exp\left[-\frac{m_1|\bm v_1-\bm U|^2+
                       m_2|\bm v_2-\bm U|^2}{2k_{\rm B}T}\right],
\]
where \(T\) is their common thermodynamic temperature and \(k_{\rm B}\) is
Boltzmann's constant.  Each velocity variance is then \(k_{\rm B}T/m_i\).
This equilibrium applies when the masses are those of the colliding objects.
Using unequal marker weights to represent identical microscopic particles
requires a separately justified sampling and collision rule.

In the equal-mass test, the velocity-
histogram residual should decrease as \(N_p^{-1/2}\), while a two-stream state
should relax toward isotropy without violating pair momentum or kinetic
energy.  For a constant cross section, compare the sampled event rate with
the analytic Maxwellian rate and compare the sampled clock in equations
(PIC30)--(PIC30a) with the collision time inferred from the same distribution.
The NTC routine continues to subcycle with its strict pair majorant after an
outer collision event begins.

Particle deposition has sampling noise.  VENUS splits markers into two
independent identifier batches, deposits each batch with the correct mass
normalization, and records their differences in density, radial current,
kinetic energy, and radial energy flux.  These differences provide a measured
noise scale for the current marker set.

Increase the marker count by factors of two and
check density, radial current, velocity dispersion, energy drift, and radial
energy flux.  Also refine the radial mesh, phase space reconstruction mesh,
ordinary timestep, potential and collision clocks, collision-shell
population, and any population-control threshold.  Repeat stochastic
comparisons with independent seeds and require deterministic resolution
differences to be smaller than the measured Monte--Carlo envelope.  The
heat-flux comparison is particularly important for gravothermal evolution
because it is more sensitive to marker noise than density.  A
zero-cross-section evolution with the same initial markers isolates
collisionless deposition and force noise from physical collisional change.
